\documentclass[aps,prd,onecolumn,amsmath,amssymb]{revtex4-2}

\usepackage{amsmath}
\usepackage{amssymb}
\usepackage{amsthm}
\usepackage{graphicx}
\usepackage{hyperref}
\usepackage{xcolor}
\usepackage{booktabs}
\usepackage{aas_macros}

\begin{document}

\title{Siegert Relation Violation as a Probe of Spacetime Foam}

\author{A.\ G.\ Kim}
\email{agkim@lbl.gov}
\affiliation{Physics Division, Lawrence Berkeley National Laboratory, Berkeley, California 94720, USA}
\author{P.\ E.\ Nugent}
\affiliation{Applied Mathematics and Computational Research Division, Lawrence Berkeley National Laboratory, Berkeley, California 94720, USA}
\author{L.\ Wang}
\affiliation{Department of Physics and Astronomy, 4242 TAMU, College Station, Texas 77843-4242, USA}

\date{\today}

\begin{abstract}
Quantum gravity predicts that spacetime fluctuates stochastically at the Planck
scale --- a ``foam'' that imprints a tiny achromatic arrival-time jitter on light
from distant sources.
We show that this jitter affects first- and second-order temporal coherence,
$g^{(1)}$ and $g^{(2)}$, differently, producing a violation of the Siegert relation
$g^{(2)}(\tau) - 1 = |g^{(1)}(\tau)|^2$ that otherwise holds for all Gaussian chaotic
sources, when the foam fluctuates on timescales shorter than the optical coherence time.
For an extended source the Siegert discrepancy is attenuated by an
angular-size factor $\Phi_{\rm eff}(r_\perp)\in[0,1]$, so the
point-source signal is the upper bound; the angular-size dependence
itself provides a differential probe of the foam transverse coherence
length $r_\perp$.
Strongly lensed sources --- whether persistent or transient --- combine
cosmological path length for maximum foam-variance accumulation with
magnification-boosted photon rates, and represent the highest-priority
targets when accessible.
These observables provide constraints that are statistically competitive with and complementary to existing imaging tests, though exquisite control of systematic errors is required to access the full range of physically viable foam models.
\end{abstract}

\maketitle

\section{Introduction}
\label{sec:intro}

A long-standing expectation from the union of quantum mechanics and general
relativity is that the smooth pseudo-Riemannian manifold of classical spacetime
breaks down at the Planck length
$\ell_{\rm P} = \sqrt{\hbar G/c^3} \simeq 1.616\times 10^{-35}$~m.
In Wheeler's original ``spacetime foam'' picture~\cite{Wheeler1955}, quantum
fluctuations of the metric on this scale render the geometry stochastic;
similar conclusions are reached in the Euclidean path-integral
treatment~\cite{Hawking1978} and in modern holographic and entropic-gravity
arguments from very different starting points.
While a complete theory of quantum gravity remains elusive, the question of
whether residual imprints of Planck-scale geometric fluctuations can accumulate
over cosmological baselines and become observable is one of the few empirical
handles on physics in the Planck regime, and it has motivated a substantial
body of phenomenological work over the past two
decades~\cite{Carlip2023,NgPerlman2022}.

A widely studied parametrization due to Ng and
collaborators~\cite{Ng2001,Ng2003,NgPerlman2022}, the $\alpha$-foam model,
encodes the foam dynamics in a power-law path-length variance
\begin{equation}
  \sigma_\ell^2(D) = A_\alpha\,D^{2(1-\alpha)}\,\ell_{\rm P}^{2\alpha},
  \label{eq:NgPerlmanVar}
\end{equation}
where $D$ is the proper distance to the source, $\alpha$ parametrizes the foam
model, and $A_\alpha$ is a dimensionless coefficient of order unity.
A photon traversing distance $D$ passes through approximately
$N \sim D/\ell_{\rm P}$ independent Planck-scale spacetime cells, each imparting
a tiny random displacement to the photon path; the exponent $\alpha$ controls
how the per-cell displacements accumulate: for random-walk foam ($\alpha=1/2$),
contributions add in quadrature ($\sigma_\ell \propto \sqrt{D\ell_{\rm P}}$);
holographic foam ($\alpha=2/3$) grows more slowly, constrained by an area-law
bound on the number of independent degrees of freedom per Planck area; and
conservative foam ($\alpha=1$) yields a distance-independent
$\sigma_\ell \propto \ell_{\rm P}$, corresponding to perfectly correlated cell
contributions.

The primary observational probe of stochastic foam to date has been the imaging
channel. Spacetime foam accumulates a cumulative phase dispersion
$\Delta\phi \simeq 2\pi\ell^{1-\alpha}\ell_{\rm P}^{\alpha}/\lambda$ across the
entrance aperture, degrading PSF sharpness and erasing interference fringes once
$\Delta\phi \gtrsim 2\pi$~\cite{LieuHillman2003}.
For PKS\,1413$+$135 ($\ell = 1.216$~Gpc) observed with HST at $\lambda=1.6\,\mu$m,
the random-walk model ($\alpha=1/2$) predicts $\Delta\phi \sim 10\times 2\pi$,
so the detection of Airy rings in this image marginally rules out that model;
the holographic model ($\alpha=2/3$) predicts $\Delta\phi \sim 10^{-9}\times 2\pi$
and is obviously not ruled out~\cite{NgPerlman2022,Christiansen2011,Perlman2015}.
Models with long-range transverse correlations evade even the random-walk
bound~\cite{LeeZurekChen2024}.

We develop a framework based on the first- and second-order coherence
functions $g^{(1)}(\tau)$ and $g^{(2)}(\tau)$ of light from astrophysical
sources.
The first-order coherence -- the field--field autocorrelation, equivalently
the fringe visibility at delay $\tau$ -- inherits a foam suppression that
interpolates continuously between unity and the Lieu--Hillman fringe-erasure
limit~\cite{LieuHillman2003}, recovering the Strehl and fringe-visibility
tests as limiting cases of the same statistic.
The second-order coherence -- the intensity--intensity autocorrelation
measured in a Hanbury Brown--Twiss experiment~\cite{HBT1956} -- is to leading
order blind to the carrier-phase modulation that suppresses $g^{(1)}$, but
acquires a distinct signature from the differential timing jitter between
photon arrivals.
The two together violate the Siegert relation
$g^{(2)}(\tau) - 1 = |g^{(1)}(\tau)|^2$~\cite{Siegert1943} in a foam-specific
way, providing a discriminant that does not require the intrinsic source
visibility to be modeled.
The finite-lag Siegert violation requires the foam to fluctuate on timescales
shorter than the optical coherence time; in the quasi-static limit
$\tau_{\rm foam}\to\infty$ the signal vanishes.  For an extended source the
Siegert discrepancy is further attenuated by an angular-size factor
$\Phi_{\rm eff}(r_\perp)\in[0,1]$, with the point-source limit
($\Phi_{\rm eff}=1$) saturating the upper bound; the source-size
dependence is itself a discriminant for the foam transverse coherence
length $r_\perp$.

Intuitively, $g^{(1)}$ measures the phase modulation of the electric field by the
foam. The second-order coherence $g^{(2)}$ measures the
differential arrival-time jitter between two photons. If both photons traverse
identical foam fluctuations, their relative timing is unchanged and $g^{(2)}$ is
unaffected. A change in $g^{(2)}$ requires the two photons to experience
\emph{different} foam: typically they follow the same path but are separated
in time by more than the foam coherence time $\tau_{\rm foam}$, so the foam
has changed between their passages.  This temporal mechanism drives the
finite-lag Siegert violation; at zero lag two photons from the same sky
position experience identical foam (their relative timing is unchanged),
and the single-mode bunching identity $g^{(2)}(0) = 2$ is preserved
exactly even for extended sources (Section~\ref{sec:extended}).

The framework draws on a well-established phenomenology of phase screens
acting on chaotic light. For a thermal point source viewed through a
\emph{static} phase screen, $g^{(1)}$ is suppressed by the Strehl factor
$e^{-\sigma_\phi^2}$ while $g^{(2)}$ is unchanged --- the principle behind
Hanbury Brown \& Twiss's demonstration that intensity interferometry
survives atmospheric phase distortions that erase Michelson
fringes~\cite{HBT1956}. A \emph{dynamic} phase screen that fluctuates
within the detection timing window imprints additional intensity
correlations on $g^{(2)}(\tau)$, the principle exploited in
Martienssen \& Spiller's pseudo-thermal light source~\cite{MartienssenSpiller1964}
and a recurring theme in HBT-through-turbulence
studies~\cite{Goodman2007,MandelWolf1995}. Planck-scale foam acts as a
dynamic phase screen distributed along the entire line of sight: when
$\tau_{\rm foam}\lesssim\tau_L$ it produces a finite-lag Siegert
discrepancy on point-like targets, attenuated by $\Phi_{\rm eff}$ for
sources whose angular extent resolves the foam transverse coherence
scale $r_\perp$.  The present work extends these laboratory phenomena
--- static-screen survival and dynamic-screen imprint --- to
cosmological propagation through spacetime foam.

Next-generation sensitivity is achievable via two complementary paths:
large-aperture single telescopes measuring the HBT effect, and multi-telescope
intensity interferometer arrays such as
VERITAS~\cite{Abeysekara2020,Acharyya2024} --- both increase collecting area
and improve signal-to-noise for nearby sources, while the temporal
$g^{(1)}$ channel is accessible to high-resolution spectroscopy.
A detection campaign across targets of varying angular size would
jointly probe $\{\alpha, r_\perp\}$ through the angular-size dependence
of the Siegert discrepancy, providing a cross-check against common
instrumental systematics.

The remainder of the paper is organized as follows.
Section~\ref{sec:model} develops the model: the spatiotemporal covariance of
the foam-induced path-length fluctuation field, the differential-jitter
distribution at a single ray, and its generalization to ray pairs.
Section~\ref{sec:g1} develops the full coherence framework for point sources:
first-order coherence and its closed-form results for Lorentzian and Voigt
line profiles; second-order coherence and its narrow-band insensitivity to
carrier phase; the resulting Siegert violation and its canonical detection
observable; and the connection to the published Strehl and fringe-visibility
tests.
Section~\ref{sec:extended} extends the analysis to spatially resolved
sources, derives the angular-size suppression of the Siegert discrepancy
by $\Phi_{\rm eff}$, shows that no foam-induced zero-lag bunching
suppression arises in this model, and identifies distance scaling and
achromaticity as discriminants against conventional propagation phase
screens.
Section~\ref{sec:prospects} assesses detectability via single-aperture HBT
measurements, derives aperture and signal-to-noise scaling, and surveys
candidate astrophysical targets.
Section~\ref{sec:conclusions} summarizes the main results.
The corrected $g^{(2)}_{\rm obs}$ derivation underlying §\ref{sec:extended}
is given in Appendix~\ref{app:g2ext}.

% ==============================================================
\section{Spacetime Foam Model}
\label{sec:model}
% ==============================================================

\subsection{Arrival-Time Jitter as a Time Reparameterization}

We model the detected field as a stochastically time-warped version of the
intrinsic source field:
\begin{equation}
  E_{\rm obs}(t) = E_{\rm src}\!\left(t - \delta t(t)\right),
  \label{eq:Eobs}
\end{equation}
where $E_{\rm src}(t)$ is a single spatial mode of a stationary zero-mean
complex Gaussian process (chaotic light, appropriate for virtually all known
astrophysical sources) and $\delta t(t)$ is the foam-induced
arrival-time jitter, a stationary stochastic process statistically independent
of $E_{\rm src}$.
The jitter arises from cumulative path-length fluctuations along the
line of sight of propagation distance $D$:
\begin{equation}
  \delta t(t) = \frac{1}{c}\,\Delta\ell_{\rm tot}(t).
  \label{eq:jitter}
\end{equation}
In the geometric-optics limit the path-length shift $\Delta\ell_{\rm tot}$
is a purely geometric quantity and hence frequency-independent (achromatic),
in contrast to plasma dispersion which scales as $\lambda^2$.
This achromaticity is a key discriminant between foam and conventional
propagation effects.

\subsection{Spatial and Temporal Fluctuations of Spacetime Foam}
\label{sec:Calpha_origin}

The scalar variance $\sigma_\ell^2(D)$ (Eq.~\eqref{eq:NgPerlmanVar})
characterizes the amplitude of foam-induced path-length fluctuations along
a single ray; here we describe how those fluctuations are correlated across
ray pairs (spatial coherence) and across time (temporal decorrelation),
both of which enter the HBT observables (see also Ref.~\cite{Carlip2023}
for a review of foam models and observational constraints).

Because foam is a property of spacetime rather than of any preferred direction,
the same dimensional logic that governs longitudinal path-length fluctuations
motivates the existence of a characteristic transverse coherence length $r_\perp$.
However, the longitudinal and transverse cases are geometrically distinct: in the
longitudinal case $D$ counts the number of cells traversed along the ray, while in
the transverse case $b$ is a separation between rays that share cells along the line
of sight --- the number of shared cells depends on both $b$ and $D$ in a way that is
not captured by a simple dimensional analogy.
Accordingly, $r_\perp$ is treated as an independent free parameter in
Section~\ref{sec:extended}.
We denote the normalized transverse coherence function by $\Phi(b/r_\perp)$,
satisfying $\Phi(0)=1$ and $\Phi(b/r_\perp)\to 0$ for $b\gg r_\perp$.
The functional form of $\Phi$ depends on the specific foam model.
For the pixellon model, Lee, Zurek \& Chen~\cite{LeeZurekChen2024} find that
the transverse correlation structure involves long-range angular correlations
(a Legendre polynomial series in aperture-plane angle) that do not reduce to
a simple $\Phi(b/r_\perp)$ ansatz, but whose long-range character strongly
suppresses image blurring relative to the uncorrelated case.
In the present framework $\Phi$ is treated as a free function; a foam
signal would manifest as an angular-size dependence of
$\mathcal{S}(\tau)$ through the suppression factor $\Phi_{\rm eff}$
(Section~\ref{sec:extended}).

Turning to the temporal evolution of spacetime foam --- an aspect not specified by Eq.~\eqref{eq:NgPerlmanVar} --- quantum foam models posit that each Planck-scale cell along the line of sight fluctuates in time.
Two photons traversing the same cell at coordinate times separated by $\tau$
experience path-length fluctuations that are correlated for $\tau \ll \tau_{\rm foam}$
but decorrelated for $\tau \gg \tau_{\rm foam}$, where $\tau_{\rm foam}$ is the
characteristic temporal fluctuation timescale at each cell.
We introduce $C_\alpha(\tau;D)$ as the normalized \emph{temporal} autocorrelation
of the per-cell path-length fluctuation, integrated over the $\sim D/\ell_{\rm P}$
cells along the line of sight.
The timescale $\tau_{\rm foam}$ is an additional free parameter of the foam model,
not fixed by the $\alpha$-model variance relation alone.
If $\tau_{\rm foam}$ is of order the Planck time,
$\ell_{\rm P}/c \approx 5\times10^{-44}$~s, then $C_\alpha(\tau) \approx 0$ for all
accessible lags: the foam is fully decorrelated on any observable timescale, the most
optimistic scenario for detection.
In the opposite \emph{quasi-static} limit $\tau_{\rm foam} \gg D/c$ (foam coherent
across the entire line of sight), $C_\alpha(\tau) \approx 1$ for all $\tau$ of
interest and every $C_\alpha$-dependent temporal observable vanishes:
the
finite-lag Siegert discrepancy $\mathcal{S}(\tau)$ is suppressed by
$[1-C_\alpha(\tau)] \to 0$ (Eq.~\eqref{eq:Sexplicit}), and the
single-ray $g^{(1)}_{\rm obs}$ recovers its intrinsic value because the
foam suppression factor
$\mathcal{F}_{\rm foam} = e^{-\sigma_\phi^2[1-C_\alpha]} \to 1$
(Eq.~\eqref{eq:Ffoam}).  In this quasi-static limit the HBT framework
loses sensitivity altogether and the foam manifests only through the
imaging-channel $g^{(1)}$ suppression on time-averaged exposures.
The full spatiotemporal covariance between two rays separated by transverse
distance $b$ takes the factorized form
\begin{equation}
  \big\langle \Delta\ell_{\rm tot}(t;\boldsymbol{\theta})\,
              \Delta\ell_{\rm tot}(t+\tau;\boldsymbol{\theta}')\big\rangle
  = \sigma_\ell^2(D)\,C_\alpha(\tau;D)\,
    \Phi\!\left(\frac{b(\boldsymbol{\theta},\boldsymbol{\theta}')}{r_\perp}\right),
  \label{eq:pathcov}
\end{equation}
where $b(\boldsymbol{\theta},\boldsymbol{\theta}') \simeq D|\boldsymbol{\theta}-\boldsymbol{\theta}'|$.
The single-ray temporal covariance appropriate for a
point source is the special case
$\boldsymbol{\theta}=\boldsymbol{\theta}'$ (so $b=0$):
\begin{equation}
  \big\langle \Delta\ell_{\rm tot}(t)\,
              \Delta\ell_{\rm tot}(t+\tau) \big\rangle
  = \sigma_\ell^2(D)\,C_\alpha(\tau;D).
  \label{eq:pathcov1ray}
\end{equation}
This product form --- separating the total variance from the normalized
spatiotemporal shape --- is a structural assumption, not a consequence of
Eq.~\eqref{eq:NgPerlmanVar}, that defines a tractable class of foam models
within which the analytic framework is exact.
The parameters $A_\alpha$, $C_\alpha(\tau;D)$, $\tau_{\rm foam}$, $r_\perp$, and
$\Phi_{\rm eff}$ parameterize distinct physical effects --- single-ray variance,
temporal foam dynamics, transverse coherence, and source angular structure ---
and each governs a distinct aspect of the signal's detectability.
The factorization is consistent with phenomenologies in which the
longitudinal path-length variance, the temporal decorrelation, and the
transverse coherence arise from statistically independent
mechanisms --- e.g.\ the simple cellular pictures of
Refs.~\cite{Ng2001,Ng2003,NgPerlman2022} where each Planck-scale cell
contributes an independent random kick.

The factorization is \emph{not} generic: the
pixellon model of Lee, Zurek \& Chen~\cite{LeeZurekChen2024} has
long-range angular correlations whose spatiotemporal covariance does
not reduce cleanly to $\sigma_\ell^2(D)\,C_\alpha(\tau)\,\Phi(b/r_\perp)$;
the framework here applies to that model only as a heuristic
approximation in the regime where transverse correlations have
decoupled.

Observationally, factorization predicts that the extended-source Siegert
discrepancy is suppressed relative to its point-source value by the
angular-size factor $\Phi_{\rm eff}(r_\perp)$:
$\mathcal{S}_{\rm ext}(\tau) \approx \Phi_{\rm eff}\,\mathcal{S}_{\rm pt}(\tau)$
in the weak-foam regime (Section~\ref{sec:extended}), independent of both
$\tau$ and $D$.  A dependence of this ratio on lag or source redshift
would falsify the factorization ansatz and indicate that a non-factorizing
model (e.g.\ pixellon) is required.
The single observable developed below --- the Siegert discrepancy
$\mathcal{S}(\tau)$ --- descends from Eq.~\eqref{eq:NgPerlmanVar}, with
Sections~\ref{sec:g1}--\ref{sec:siegert} probing its longitudinal
projection (single-ray temporal jitter) and Section~\ref{sec:extended}
its transverse projection (cross-ray $\Phi_{\rm eff}$ suppression).

\subsection{Differential Jitter and Its Distribution}
\label{sec:differential_jitter}

With the spatiotemporal covariance model in place, we define the key derived quantity
that enters the coherence integrals.
Define the \emph{cross-ray differential jitter} between a photon arriving along
direction $\boldsymbol{\theta}$ at time $t+\tau$ and a photon arriving along
$\boldsymbol{\theta}'$ at time $t$:
\begin{equation}
  \Delta_{\boldsymbol{\theta},\boldsymbol{\theta}'}(\tau)
  \equiv \delta t(t+\tau;\boldsymbol{\theta}) - \delta t(t;\boldsymbol{\theta}').
  \label{eq:DeltaExt}
\end{equation}
For fixed $D$, stationarity implies the statistics depend only on $\tau$ and the
transverse separation $b(\boldsymbol{\theta},\boldsymbol{\theta}')=D|\boldsymbol{\theta}-\boldsymbol{\theta}'|$.
Expanding the square and applying Eq.~\eqref{eq:pathcov}:
\begin{align}
  \sigma_\Delta^2(\tau;\boldsymbol{\theta},\boldsymbol{\theta}')
  &= \frac{1}{c^2}\big\langle
    [\Delta\ell_{\rm tot}(t+\tau;\boldsymbol{\theta})
     - \Delta\ell_{\rm tot}(t;\boldsymbol{\theta}')]^2
    \big\rangle \notag \\
  &= \frac{1}{c^2}\left[
    \sigma_\ell^2 + \sigma_\ell^2
    - 2\sigma_\ell^2\,C_\alpha(\tau)\,\Phi\!\left(\frac{b}{r_\perp}\right)
    \right] \notag \\[4pt]
  &= \frac{2\sigma_\ell^2(D)}{c^2}
    \left[1 - C_\alpha(\tau;D)\,\Phi\!\left(\frac{b(\boldsymbol{\theta},\boldsymbol{\theta}')}{r_\perp}\right)\right].
  \label{eq:diffVarExt}
\end{align}
For $b\gg r_\perp$, $\Phi\to 0$ and the variance becomes
$2\sigma_\ell^2/c^2$ \emph{at all lags, including $\tau=0$} — a new effect
absent in the point-source case, and one that is independent of the temporal
fluctuation structure of the foam (i.e.\ independent of $C_\alpha$).
In the dynamical limit $C_\alpha\to 0$ this result holds for all separations
$b$, independently of $\Phi$.

By the central limit theorem (CLT), $\Delta_{\boldsymbol{\theta},\boldsymbol{\theta}'}(\tau)$
is a sum of contributions from many foam cells and is therefore Gaussian with zero mean: 
\begin{equation}
  \Delta_{\boldsymbol{\theta},\boldsymbol{\theta}'}(\tau)
  \sim \mathcal{N}\!\left(0,\,\sigma_\Delta^2(\tau;\boldsymbol{\theta},\boldsymbol{\theta}')\right),
  \label{eq:kernel}
\end{equation}
where $\sigma_\Delta^2$ is taken from Eq.~\eqref{eq:diffVarExt}.

The Gaussian approximation follows from the CLT: for any source at
$D \gtrsim 1$~kpc, the cell count $N_{\rm cells} \sim D/\ell_{\rm P} \gtrsim 10^{46}$
amply satisfies the large-$N$ requirement and higher-cumulant corrections
are negligible.
For the cross-ray case the additional condition $b \gg r_\perp$ is assumed,
and is satisfied in the fully-resolved limit $\Phi_{\rm eff} \to 0$ adopted
in Section~\ref{sec:prospects}.

% ==============================================================
\section{Coherence Functions and Siegert Violation}
\label{sec:g1}
% ==============================================================

\subsection{Point-Source Differential Jitter}
\label{sec:pointsource_jitter}

In this section we restrict to a point source ($b=0$),
for which the cross-ray differential jitter of Eq.~\eqref{eq:DeltaExt}
reduces to the single-ray differential jitter
\begin{equation}
  \Delta(\tau) \equiv \delta t(t+\tau) - \delta t(t),
  \label{eq:DeltaDef}
\end{equation}
with variance obtained from Eq.~\eqref{eq:diffVarExt} at $\Phi(0)=1$:
\begin{align}
  \sigma_\Delta^2(\tau;D)
  &= \frac{2\sigma_\ell^2(D)}{c^2}\big[1 - C_\alpha(\tau;D)\big].
  \label{eq:diffVar}
\end{align}

Although this section restricts to a point source for notational compactness, this single-ray variance applies equally to extended sources at finite lag $\tau \neq 0$: under the factorized covariance Eq.~\eqref{eq:pathcov}, the temporal factor $C_\alpha(\tau;D)$ is identical for all ray pairs, so cross-ray transverse structure enters only at $\tau=0$ through $\Phi_{\rm eff}$ (Section~\ref{sec:extended}).

\subsection{Jitter Convolution Formula for $g^{(1)}$}

The observed first-order coherence is
\begin{equation}
  g^{(1)}_{\rm obs}(\tau) = \frac{\langle E^*_{\rm obs}(t)\,E_{\rm obs}(t+\tau)\rangle}
                                  {\langle|E_{\rm obs}(t)|^2\rangle}.
  \label{eq:g1def}
\end{equation}
Substituting Eq.~\eqref{eq:Eobs} and letting
$u = t - \delta t(t)$, $v = t+\tau - \delta t(t+\tau)$, so that
$v - u = \tau - \Delta(\tau)$, stationarity of $E_{\rm src}$ gives
\begin{equation}
  \langle E^*_{\rm src}(u)\,E_{\rm src}(v)\rangle
  = \langle|E_{\rm src}|^2\rangle\,g^{(1)}_{\rm src}(\tau - \Delta(\tau)).
\end{equation}
Since $\delta t$ is independent of $E_{\rm src}$, averaging over the jitter
yields
\begin{equation}
  g^{(1)}_{\rm obs}(\tau)
    = \big\langle g^{(1)}_{\rm src}(\tau - \Delta(\tau))\big\rangle_{\delta t}
    = \int_{-\infty}^{\infty} \mathcal{N}(\Delta;0,\sigma_\Delta^2(\tau))\,
      g^{(1)}_{\rm src}(\tau - \Delta)\,d\Delta.
  \label{eq:g1obs}
\end{equation}
The observed first-order coherence is the intrinsic coherence function
convolved with the Gaussian differential-jitter kernel.

\subsection{Closed Forms for Lorentzian and Voigt Sources}

As a representative case admitting a closed-form result, we take a Lorentzian
line profile, which corresponds to an exponential coherence function via the
Wiener--Khinchin theorem:
\begin{equation}
  g^{(1)}_{\rm src}(\tau) = e^{-|\tau|/\tau_L},
  \label{eq:g1src_lorentz}
\end{equation}
where $\tau_L$ is the coherence time.
Evaluating Eq.~\eqref{eq:g1obs} by completing the square and splitting on the
sign of $\tau-\Delta$ gives, for $\tau > 0$:
\begin{equation}
  g^{(1)}_{\rm obs}(\tau)
  = e^{\sigma_\Delta^2/2\tau_L^2}
    \cdot\frac{1}{2}\left[
      e^{-\tau/\tau_L}\,\mathrm{erfc}\!\left(\frac{\sigma_\Delta/\tau_L - \tau/\sigma_\Delta}{\sqrt{2}}\right)
      + e^{\tau/\tau_L}\,\mathrm{erfc}\!\left(\frac{\sigma_\Delta/\tau_L + \tau/\sigma_\Delta}{\sqrt{2}}\right)
    \right].
  \label{eq:g1obs_erfc}
\end{equation}
In the weak-jitter limit $\sigma_\Delta \ll \tau_L$, the Gaussian average of
$e^{-|\tau-\Delta|/\tau_L}$ over $\Delta\sim\mathcal{N}(0,\sigma_\Delta^2)$
is dominated by $|\Delta|\ll|\tau|$, so one expands around $\Delta=0$:
\begin{equation}
  g^{(1)}_{\rm obs}(\tau) \approx e^{-|\tau|/\tau_L}
  \left[1 + \frac{\sigma_\Delta^2(\tau)}{2\tau_L^2} + \mathcal{O}(\sigma_\Delta^4)\right].
  \label{eq:g1weak}
\end{equation}
Here $\sigma_\Delta(\tau)$ is $\tau$-dependent through $C_\alpha(\tau)$,
and the expansion is in $\sigma_\Delta(\tau)/\tau_L$ at fixed $\tau$.
The correction is positive because $e^{-|\tau|/\tau_L}$ is convex and
Jensen's inequality applies.
The normalization $g^{(1)}_{\rm obs}(0) = 1$ is preserved exactly.

For a Voigt profile $g^{(1)}_{\rm src}(\tau) = e^{-|\tau|/\tau_L}\,e^{-\tau^2/2\tau_G^2}$
--- a Lorentzian of width $\tau_L$ (from optically-thick electron scattering or
BLR turbulent broadening) convolved with a Gaussian of width $\tau_G$ (from
bulk Doppler motions) --- the jitter kernel combines with the source Gaussian
into an effective width and mean,
\begin{equation}
  \tilde\sigma_1^2 = \frac{\sigma_\Delta^2\,\tau_G^2}{\sigma_\Delta^2+\tau_G^2},
  \qquad
  \mu_1(\tau) = \frac{\tau_G^2\,\tau}{\sigma_\Delta^2+\tau_G^2},
  \label{eq:voigt_params}
\end{equation}
giving, for $\tau > 0$:
\begin{equation}
  g^{(1)}_{\rm obs}(\tau)
  = \frac{\tilde\sigma_1}{\sigma_\Delta}\,
    e^{-\tau^2/2(\sigma_\Delta^2+\tau_G^2)}\,
    e^{\tilde\sigma_1^2/2\tau_L^2}\,
    \frac{1}{2}\!\left[
      e^{-\mu_1/\tau_L}\,\mathrm{erfc}\!\left(\frac{\tilde\sigma_1/\tau_L - \mu_1/\tilde\sigma_1}{\sqrt{2}}\right)
      +e^{\mu_1/\tau_L}\,\mathrm{erfc}\!\left(\frac{\tilde\sigma_1/\tau_L + \mu_1/\tilde\sigma_1}{\sqrt{2}}\right)
    \right].
  \label{eq:g1obs_voigt}
\end{equation}
The normalization $g^{(1)}_{\rm obs}(0)=1$ is preserved: recall
$\sigma_\Delta = \sigma_\Delta(\tau)$ with $\sigma_\Delta(0)=0$
(Eq.~\eqref{eq:diffVar}), so $\tilde\sigma_1(0)=0$, the exponential prefactors
approach unity, and the erfc terms reduce to
$\mathrm{erfc}(0)+\mathrm{erfc}(0)=2$.
In the limit $\tau_G\to\infty$ (pure Lorentzian), $\tilde\sigma_1\to\sigma_\Delta$ and
$\mu_1\to\tau$, so Eq.~\eqref{eq:g1obs_voigt} reduces to Eq.~\eqref{eq:g1obs_erfc}.
In the limit $\tau_L\to\infty$ (pure Gaussian), the erfc terms combine to unity and
$g^{(1)}_{\rm obs}(\tau) = (\tilde\sigma_1/\sigma_\Delta)\,e^{-\tau^2/2(\sigma_\Delta^2+\tau_G^2)}$:
the observed coherence is a Gaussian whose width is broadened in quadrature by the jitter.

\subsection{Narrow-Band Limit}
\label{sec:narrowband}

For a quasi-monochromatic source with center frequency $\omega_0$ and
angular-frequency bandwidth $\Delta\omega \ll \omega_0$, 
$E_{\rm src}(t) = A(t)\,e^{-i\omega_0 t}$
where $A(t)$ varies on the coherence timescale
$\tau_c = 1/\Delta\omega = \lambda^2/(2\pi c\,\Delta\lambda)$.
When $\sigma_\ell/(c\tau_c) \ll 1$, typical realizations of $\delta t$ leave the
envelope unaffected, $A(t-\delta t) \approx A(t)$, and the jitter acts purely as
a carrier phase modulation; the observed field takes the form
\begin{equation}
  E_{\rm obs}(t) \approx A(t)\,e^{-i\omega_0 t}\,e^{i\omega_0\delta t(t)}.
  \label{eq:Eobs_narrowband}
\end{equation}
(Corrections to the envelope approximation are of order $(\sigma_\ell/c\tau_c)^2$;
for $\alpha=1/2$ at $D=1$~Gpc, $\lambda=500$~nm, $\Delta\lambda/\lambda=10^{-3}$
this parameter is $\approx 0.28$, so the exact convolution Eq.~\eqref{eq:g1obs_erfc}
should be used instead.
For $\alpha=2/3$ the analogous parameter is ${\sim}10^{-11}$
(Section~\ref{sec:scales}), so Eq.~\eqref{eq:g1narrowband} is exact to any
practical precision in the holographic case.)
Since $A(t)$ and $\delta t(t)$ are independent, the first-order coherence
factorizes.
The jitter contribution is the characteristic function of $\Delta(\tau)$
evaluated at $\omega_0$; since $\Delta(\tau)$ is Gaussian
(Section~\ref{sec:differential_jitter}),
\begin{equation}
  \big\langle e^{i\omega_0\Delta(\tau)}\big\rangle_{\delta t}
  = e^{-\omega_0^2\sigma_\Delta^2(\tau)/2}.
\end{equation}
This holds when $\tau \gg \tau_{\rm foam}$ so that foam-cell contributions are
independent; in the dynamical limit $\tau_{\rm foam} \ll \tau_c$ the condition
is satisfied at all observationally relevant lags.
Substituting $\omega_0 = 2\pi c/\lambda$ and Eq.~\eqref{eq:diffVar}:
\begin{equation}
  g^{(1)}_{\rm obs}(\tau)
    = g^{(1)}_{\rm src}(\tau)\cdot
      \exp\!\left(-\sigma_\phi^2\big[1 - C_\alpha(\tau;D)\big]\right),
  \label{eq:g1narrowband}
\end{equation}
where we define the rms accumulated phase
\begin{equation}
  \sigma_\phi^2 \equiv \left(\frac{2\pi}{\lambda_{\rm obs}}\right)^2
  \int_0^{L} \big[1+z(s)\big]^2\,\frac{d\sigma_\ell^2}{ds}\,ds,
  \label{eq:sigmaphi}
\end{equation}
integrating the foam differential variance
$d\sigma_\ell^2/ds = A_\alpha\,2(1-\alpha)\,s^{1-2\alpha}\,\ell_{\rm P}^{2\alpha}$
along the photon's local proper path of total length $L$. The $(1+z)^2$
weighting reflects that a foam cell at redshift $z$ imparts an observer-frame
time delay $\delta t_{\rm obs} = (1+z)\,\delta\ell_{\rm local}/c$, and hence
a phase $\delta\phi = (2\pi/\lambda_{\rm obs})\,(1+z)\,\delta\ell_{\rm local}$,
because the photon's wavelength at $z$ is
$\lambda_{\rm local}(z) = \lambda_{\rm obs}/(1+z)$.
Switching integration variable to comoving distance via
$ds = dr_{\rm C}/(1+z)$ yields the equivalent form
$\sigma_\phi^2 = (2\pi/\lambda_{\rm obs})^2 \int_0^{D_{\rm C}}
(1+z)\,(d\sigma_\ell^2/ds)\,dr_{\rm C}$.
For $z \ll 1$ both forms reduce to the flat-space limit
$\sigma_\phi = 2\pi\sqrt{A_\alpha}\,D^{1-\alpha}\,\ell_{\rm P}^\alpha/\lambda$.
We write the foam suppression factor as
\begin{equation}
  \mathcal{F}_{\rm foam}(\tau) \equiv e^{-\sigma_\phi^2[1-C_\alpha(\tau;D)]},
  \label{eq:Ffoam}
\end{equation}
so $g^{(1)}_{\rm obs}(\tau) = g^{(1)}_{\rm src}(\tau)\cdot\mathcal{F}_{\rm foam}(\tau)$.
At $\tau=0$: $\mathcal{F}_{\rm foam}(0) = 1$.
At $\tau\to\infty$: $\mathcal{F}_{\rm foam}(\infty) = e^{-\sigma_\phi^2[1-C_\alpha(\infty;D)]}$,
which equals $e^{-\sigma_\phi^2}$ in the dynamical limit ($C_\alpha\to 0$)
and approaches unity for quasi-static foam ($C_\alpha\approx 1$).
The narrow-band condition $\sigma_\ell/(c\tau_c)\ll 1$ is compatible with
the weak-jitter expansion of Eq.~\eqref{eq:g1weak} --- the former controls
envelope displacement, the latter the carrier-phase suppression --- and
when neither applies the full convolution Eq.~\eqref{eq:g1obs} must be
evaluated.

% ==============================================================
\subsection{Second-Order Coherence}
\label{sec:g2}
% ==============================================================
We derive the second-order coherence $g^{(2)}_{\rm obs}(\tau)$ for a point source,
continuing the framework of Section~\ref{sec:g1}.
The result has the same structure as Eq.~\eqref{eq:g1obs}: $g^{(2)}_{\rm obs}-1$
is a convolution of $g^{(2)}_{\rm src}-1$ with the identical Gaussian jitter kernel.
Because the intensity $|E_{\rm obs}|^2$ is blind to the carrier phase
$e^{i\omega_0\delta t}$, foam leaves $g^{(2)}$ nearly unchanged in the narrow-band
limit while strongly suppressing $g^{(1)}$ --- the asymmetry that produces the
Siegert violation developed in Section~\ref{sec:siegert}.

The same argument makes $g^{(2)}$ practically immune to atmospheric seeing under
current observing conditions, in contrast to $g^{(1)}$. Atmospheric turbulence
introduces path-length changes of order $\Delta\ell_{\rm atm} \sim 0.1$--$1\,\rm mm$,
far exceeding the optical wavelength; this randomizes the carrier phase and
destroys fringe visibility without adaptive optics. However, $g^{(2)}$ is blind
to this carrier-phase modulation, and the atmospheric coherence time
$\tau_0 \sim 1\,\rm ms$ far exceeds the HBT correlation window $\tau_L \sim \rm ns$,
so both photons in a coincidence traverse a frozen atmosphere and the differential
atmospheric jitter is negligible. This advantage narrows as detector timing
resolution approaches the atmospheric timing jitter
$\Delta\ell_{\rm atm}/c \sim 1$--$3\,\rm ps$; sub-ps detectors will require
atmospheric monitoring or site-selection to control this systematic.

The Wick-contraction derivation of Appendix~\ref{app:g2ext}, specialized
to a point source ($\Phi = 1$, $X = Y \equiv \Delta(\tau)$), yields
\begin{equation}
  g^{(2)}_{\rm obs}(\tau) - 1
    = \int_{-\infty}^{\infty} \mathcal{N}(\Delta;0,\sigma_\Delta^2(\tau))\,
      \big[g^{(2)}_{\rm src}(\tau-\Delta)-1\big]\,d\Delta,
  \label{eq:g2obs}
\end{equation}
with $\sigma_\Delta^2(\tau) = (2\sigma_\ell^2(D)/c^2)[1-C_\alpha(\tau;D)]$
the same single-ray differential variance that enters Eq.~\eqref{eq:g1obs}.
The structure is the same as the $g^{(1)}$ result: a Gaussian jitter
kernel convolves the intrinsic source coherence. The Gaussian kernel
requires $\tau \gg \tau_{\rm foam}$; for $\tau \lesssim \tau_{\rm foam}$
Eq.~\eqref{eq:g2obs} should be regarded as approximate
(Section~\ref{sec:Calpha_origin}).

The intensity $I_{\rm obs}(t)=|E_{\rm obs}(t)|^2$ is blind to the carrier phase
$e^{i\omega_0\delta t(t)}$: only the envelope matters.
In the narrow-band limit $\delta t \ll \tau_c$, the jitter does not appreciably
shift the envelope, $A(t-\delta t)\approx A(t)$, so
$I_{\rm obs}(t)\approx|A(t)|^2=I_{\rm src}(t)$ to leading order in $\delta t/\tau_c$.
Consequently:
\begin{equation}
  g^{(2)}_{\rm obs}(\tau) = g^{(2)}_{\rm src}(\tau)
  \qquad \text{(leading order in $\delta t/\tau_c$)},
  \label{eq:g2unchanged}
\end{equation}
with corrections of order $(\sigma_\ell/c\tau_c)^2$ from both envelope-timing
and carrier-phase contributions to the intensity.
For $\alpha=1/2$ at $D=1$~Gpc, $\lambda=500$~nm, $\Delta\lambda/\lambda=10^{-3}$:
$\delta t_{\rm rms}/\tau_c \approx 0.28$, so corrections are at the ${\sim}8\%$
level and the exact convolution Eq.~\eqref{eq:g2obs} should be used for precision
work.
For $\alpha=2/3$ the expansion parameter $\sigma_\ell/(c\tau_c)$ is $\sim 10^{-11}$ and
Eq.~\eqref{eq:g2unchanged} is exact to any practical precision.

By contrast, $g^{(1)}$ \emph{is} sensitive to the carrier phase because
$\sigma_\phi = \omega_0\sigma_\ell/c$ can be large even when $\delta t \ll \tau_c$:
for the same parameters, $\sigma_\phi \approx 280 \gg 1$.



% ==============================================================
\subsection{Siegert Relation Violation}
\label{sec:siegert}

The Siegert relation $g^{(2)}(\tau)-1=|g^{(1)}(\tau)|^2$ holds for any
stationary complex Gaussian field and provides a powerful internal consistency
check: once $g^{(1)}$ is known, the expected $g^{(2)}$ is fixed.
As foreshadowed in Section~\ref{sec:g2}, foam violates this relation because
$g^{(1)}_{\rm obs}$ is suppressed by the carrier-phase factor
$\mathcal{F}_{\rm foam}$ while $g^{(2)}_{\rm obs}$ is not.
This subsection derives and characterizes the Siegert discrepancy
$\mathcal{S}(\tau) \equiv [g^{(2)}_{\rm obs}-1] - |g^{(1)}_{\rm obs}|^2$.

\subsubsection{General Bound}

Define the Siegert discrepancy:
\begin{equation}
  \mathcal{S}(\tau)
  \equiv \left[g^{(2)}_{\rm obs}(\tau)-1\right] - \left|g^{(1)}_{\rm obs}(\tau)\right|^2.
  \label{eq:Sdef}
\end{equation}
Let $X \equiv g^{(1)}_{\rm src}(\tau-\Delta(\tau))$, a complex random variable.
From Eqs.~\eqref{eq:g1obs} and \eqref{eq:g2obs}:
\begin{equation}
  g^{(1)}_{\rm obs}(\tau) = \langle X\rangle_{\delta t}, \qquad
  g^{(2)}_{\rm obs}(\tau) - 1 = \langle |X|^2\rangle_{\delta t}.
\end{equation}
Therefore
\begin{equation}
  \mathcal{S}(\tau)
  = \langle |X|^2\rangle - |\langle X\rangle|^2
  \geq 0,
  \label{eq:Svar}
\end{equation}
by the Cauchy--Schwarz inequality, with equality only when $\Delta(\tau)$ has
zero variance (i.e., at $\tau=0$ or in the absence of foam).
The Siegert relation is violated for all $\tau\neq 0$ whenever foam is present.
The Siegert violation is not a non-classical signature: the foam-modified field
remains a classical stochastic process, and the classical inequalities
$g^{(2)}(0)\geq g^{(2)}(\tau)$ and $g^{(2)}(0)>1$ are preserved; it is only
the Gaussian-statistics identity relating $g^{(2)}$ to $g^{(1)}$ that is broken, due to the stochastic phase accumulated along the propagation path.

\subsubsection{Sensitivity Scale Decoupling}
\label{sec:scales}

The violation arises because $g^{(1)}$ and $g^{(2)}$ respond to foam on
different scales.
The first-order coherence is controlled by $\sigma_\phi =
2\pi\sigma_\ell(D)/\lambda$, the rms accumulated phase along a single ray;
$g^{(1)}$ is sensitive to it because the carrier phase $e^{i\omega_0\delta t}$
appears directly in the field correlation.
The second-order coherence is controlled by $\sigma_\Delta(\tau)/\tau_c$,
the rms differential jitter scaled by the coherence time; $g^{(2)}$ is
sensitive to it because intensity correlations depend on how much the
jitter \emph{changes} between two detection events, not its absolute value.
These two scales are related by
\begin{equation}
  \frac{\sigma_\Delta(\tau)}{\tau_c}
  = \frac{\sqrt{2[1-C_\alpha(\tau)]}\,\sigma_\ell(D)}{c\tau_c}
  = \sigma_\phi \cdot \frac{\Delta\lambda}{\lambda}\cdot\sqrt{2[1-C_\alpha(\tau)]},
  \label{eq:ratio}
\end{equation}
using $\omega_0\tau_c = \lambda/\Delta\lambda$.
The factor $\Delta\lambda/\lambda \ll 1$ means that even when $\sigma_\phi \gg 1$
(fringes erased), $\sigma_\Delta(\tau)/\tau_c$ can be small --- so for a
\emph{point source} $g^{(2)}$ remains near its intrinsic value while
$g^{(1)}$ is suppressed.  This decoupling is the physical root of the
Siegert violation.  For an extended source the decoupling is incomplete:
at finite lag $g^{(2)}_{\rm obs}$ picks up the carrier-phase scale
$\sigma_\phi$ through the cross-direction factor $\mathcal{F}_{\rm cross}$
of Eq.~\eqref{eq:Fcross}, attenuating the Siegert discrepancy by the
angular-size factor $\Phi_{\rm eff}$ (Section~\ref{sec:extended}).
At the fiducial parameters $D=1$~Gpc, $\lambda=500$~nm,
$\Delta\lambda/\lambda=10^{-3}$, $A_\alpha=1$, and the dynamical limit
$C_\alpha\to 0$, the random-walk model $\alpha=1/2$ gives
$\sigma_\phi = 280$ and $\sigma_\Delta(\tau)/\tau_c = 0.40$, while the
holographic model $\alpha=2/3$ gives $\sigma_\phi = 2.5\times 10^{-8}$ and
$\sigma_\Delta(\tau)/\tau_c = 3.6\times 10^{-11}$.
For quasi-static foam ($C_\alpha\approx 1$) the ratio
$\sigma_\Delta(\tau)/\tau_c$ vanishes at all lags.
The $\alpha=1/2$ case is excluded by imaging under the standard assumption
of short-range transverse foam correlations~\cite{Perlman2015}; models with
long-range correlations evade this bound~\cite{LeeZurekChen2024}, and
$\alpha=2/3$ (holographic) is the most sensitive surviving model.


\subsubsection{Explicit Form and Properties in the Narrow-Band Limit}

In the narrow-band limit, $g^{(2)}_{\rm obs} \approx g^{(2)}_{\rm src}$
(Eq.~\eqref{eq:g2unchanged}), while $g^{(1)}_{\rm obs}$ is suppressed by
$\mathcal{F}_{\rm foam}(\tau)$ (Eq.~\eqref{eq:g1narrowband}).
Using the intrinsic Siegert relation $g^{(2)}_{\rm src}(\tau)-1=|g^{(1)}_{\rm src}(\tau)|^2$,
to leading order:
\begin{equation}
  \mathcal{S}(\tau) \approx
  |g^{(1)}_{\rm src}(\tau)|^2\!\left(1-e^{-2\sigma_\phi^2[1-C_\alpha(\tau;D)]}\right).
  \label{eq:Sexplicit}
\end{equation}
In the weak-jitter limit $\sigma_\phi\ll 1$,
$\mathcal{S}\approx 2\sigma_\phi^2[1-C_\alpha]\,|g^{(1)}_{\rm src}|^2$ — linear
in $\sigma_\phi^2$ and small.
In the strong-jitter limit $\sigma_\phi\gg 1$, the exponential vanishes and
$\mathcal{S}\approx|g^{(1)}_{\rm src}(\tau)|^2$: $g^{(2)}_{\rm obs}-1$ retains
its intrinsic value while $|g^{(1)}_{\rm obs}|^2\to 0$, producing a full violation.
Quasi-static foam ($C_\alpha\approx 1$) gives $\mathcal{S}\approx 0$ at all lags
— foam is invisible to all coherence observables in this limit.
The discrepancy $\mathcal{S}(\tau)$ is constructed directly from
the two measurable quantities $|g^{(1)}_{\rm obs}(\tau)|$ and
$g^{(2)}_{\rm obs}(\tau)-1$ (or equivalently $\mathcal{S}(\nu)$
in the spectral domain), and the modulation beyond the prefactor $|g^{(1)}_{\rm src}(\tau)|^2$
is controlled solely by $\alpha$, $D$, and $C_\alpha(\tau;D)$, independently
of the unknown source spectrum.

For a Lorentzian source, $\mathcal{S}(\tau)$ peaks at a lag set by the
competition between the source-coherence falloff $e^{-2\tau/\tau_L}$ and
the foam-induced rise $1-e^{-2\sigma_\phi^2[1-C_\alpha(\tau)]}$.
Differentiating Eq.~\eqref{eq:Sexplicit} yields the implicit condition
\begin{equation}
  -\sigma_\phi^2\,C_\alpha'(\tau_{\max})\,e^{-2\sigma_\phi^2[1-C_\alpha(\tau_{\max})]}
  = \frac{1}{\tau_L}\left(1-e^{-2\sigma_\phi^2[1-C_\alpha(\tau_{\max})]}\right).
  \label{eq:taumax_general}
\end{equation}
In the weak-jitter limit ($\sigma_\phi^2\ll 1$) this reduces to the
$\sigma_\phi$-independent statement $-C_\alpha'(\tau_{\max})/[1-C_\alpha(\tau_{\max})]=2/\tau_L$.
For the canonical exponential decorrelation $C_\alpha(\tau)=e^{-\tau/\tau_{\rm foam}}$,
\begin{equation}
  \tau_{\max} = \tau_{\rm foam}\,\ln\!\left(1+\frac{\tau_L}{2\tau_{\rm foam}}\right),
  \label{eq:taumax_exp}
\end{equation}
which interpolates between $\tau_{\max}\approx \tau_{\rm foam}\ln[\tau_L/(2\tau_{\rm foam})]$
for fast foam ($\tau_{\rm foam}\ll\tau_L$) and $\tau_{\max}\to\tau_L/2$ for slow
foam ($\tau_{\rm foam}\gg\tau_L$).

In the fast-foam / dynamical limit
($\tau_{\rm foam}\ll\tau_L$, $C_\alpha\to 0$ above $\tau_{\rm foam}$),
$|g^{(1)}_{\rm src}(\tau_{\max})|^2\to 1$ and $1-C_\alpha(\tau_{\max})\to 1$,
so Eq.~\eqref{eq:Sexplicit} reduces to
\begin{equation}
  \mathcal{S}(\tau_{\max}) \approx 1 - e^{-2\sigma_\phi^2},
  \label{eq:Stc}
\end{equation}
the canonical detection observable: quadratic
$\mathcal{S}(\tau_{\max})\approx 2\sigma_\phi^2$ in the weak-foam limit,
saturating at $1$ for $\sigma_\phi\gg 1$. 

% ==============================================================
\subsection{Connection to Fringe Visibility and Strehl Ratio}
\label{sec:connection}
% ==============================================================

The fringe visibility and Strehl ratio are spatial-coherence observables
--- they probe phase variations across an aperture or between two
separated telescopes, described by $g^{(1)}_{\rm obs}(B)$ --- while
$g^{(1)}_{\rm obs}(\tau)$ derived above describes temporal coherence
along a single line of sight. The two cases are directly analogous: the
temporal decorrelation $C_\alpha(\tau;D)$, measuring how the foam changes
between two observations separated by time $\tau$, maps to the spatial
decorrelation $\Phi(B/r_\perp)$ (Section~\ref{sec:model}), measuring how
correlated the foam fluctuations are at two aperture points separated by
baseline $B$. In both cases the foam suppression has the same form, with
$C_\alpha$ replaced by $\Phi$:
\begin{equation}
  \mathcal{F}_{\rm foam}^{\rm spat}(B)
    \equiv e^{-\sigma_\phi^2[1-\Phi(B/r_\perp)]}.
  \label{eq:Ffoam_spat}
\end{equation}
This mirrors $\mathcal{F}_{\rm foam}(\tau)$ exactly: $\Phi(0)=1$ gives no
suppression at zero baseline, while $\Phi(B/r_\perp)\to 0$ for
$B\gg r_\perp$ gives maximum suppression $e^{-\sigma_\phi^2}$ --- just as
$\mathcal{F}_{\rm foam}(0)=1$ and
$\mathcal{F}_{\rm foam}(\infty)=e^{-\sigma_\phi^2}$.

The fringe visibility at baseline $B$ is then
\begin{equation}
  V(B) = V_{\rm src}(B)\cdot\mathcal{F}_{\rm foam}^{\rm spat}(B)
       = V_{\rm src}(B)\cdot e^{-\sigma_\phi^2[1-\Phi(B/r_\perp)]},
  \label{eq:visibility}
\end{equation}
where $V_{\rm src}(B) = |g^{(1)}_{\rm src}|$ evaluated at angular scale
$\lambda/B$. In the limit $B\gg r_\perp$ the two aperture points sample
independent foam realizations and the suppression reaches its maximum
$e^{-\sigma_\phi^2}$. Lieu \&
Hillman~\cite{LieuHillman2003} implicitly assumed this limit: their
erasure condition $\Delta\phi \geq 2\pi$ with
$\Delta\phi \simeq 2\pi\sigma_\ell(D)/\lambda$~\cite{NgPerlman2022} is
precisely $\sigma_\phi \geq 1$, the limit
$\mathcal{F}_{\rm foam}^{\rm spat}\to 0$. Equation~\eqref{eq:visibility}
extends their binary criterion to a continuous suppression curve whose
shape encodes $\Phi(B/r_\perp)$ --- in direct analogy to how
$g^{(1)}_{\rm obs}(\tau)$ encodes the temporal decorrelation
$C_\alpha(\tau;D)$.

The Strehl ratio for a single aperture of diameter $D_{\rm tel}$ follows
the same logic with $B\sim D_{\rm tel}$: when $r_\perp\ll D_{\rm tel}$,
aperture-plane phases are independent draws from the single-ray
distribution and $S=e^{-\sigma_\phi^2}$.\footnote{%
  The Mar\'{e}chal formula $S = e^{-\sigma_\phi^2}$ (valid for
  $\sigma_\phi \lesssim 1$) and the temporal result
  $\mathcal{F}_{\rm foam}(\infty) = e^{-\sigma_\phi^2}$ are linked by a
  spatio-temporal ergodic identification: both require $r_\perp \ll D_{\rm tel}$
  and $C_\alpha(\infty;D)\to 0$. For quasi-static foam ($C_\alpha\approx 1$) or
  $r_\perp \gg D_{\rm tel}$, the aperture sits within a single coherence patch,
  the wavefront carries only a common phase offset, and $S\to 1$ even when
  $\sigma_\phi^2$ is large. The transverse coherence function $\Phi(b/r_\perp)$
  governs a physically distinct effect --- cross-ray correlations at angular
  separations $\boldsymbol{\theta}\neq\boldsymbol{\theta}'$ --- which enters
  the second-order coherence of extended sources (Section~\ref{sec:extended})
  but not the single-aperture Strehl ratio.}
Both Strehl and fringe visibility therefore measure the same underlying
$\sigma_\phi^2 = (2\pi\sigma_\ell/\lambda)^2$ in the large-baseline limit
$B\gg r_\perp$, and the published
constraints~\cite{LieuHillman2003,NgPerlman2022,Perlman2015,Christiansen2011,Ragazzoni2003}
are unified within this framework.

% ==============================================================
\section{Extended Sources: Finite-Lag Suppression of the Siegert Discrepancy}
\label{sec:extended}
% ==============================================================

We now generalize from a point source to an extended incoherent source.
Three results emerge: (i) the first-order coherence $g^{(1)}_{\rm obs}$ is
unchanged from the point-source case; (ii) the second-order coherence at
zero lag is exactly $g^{(2)}_{\rm obs}(0) = 2$ --- a single-mode detector
observing chaotic light from an extended source develops no foam-induced
suppression of the zero-lag bunching peak; and (iii) at finite lag the
extended-source Siegert discrepancy is \emph{suppressed} relative to the
point-source value by an angular-size factor $\Phi_{\rm eff} \in [0,1]$,
so the point-source Siegert discrepancy is the maximum-signal upper bound
for the foam signature in this framework.  The angular-size dependence is
itself a discriminant for the foam transverse coherence length $r_\perp$
across targets of varying angular size.  The full pedestrian derivation is
given in Appendix~\ref{app:g2ext}; this section presents the key results.

\subsection{Field Model and Spatiotemporal Covariance}

We generalize to an extended incoherent source whose angular brightness
distribution is $I(\boldsymbol{\theta})$ (normalized to unity).
The analysis assumes a point (single spatial mode) detector; a finite aperture
would introduce Van~Cittert--Zernike geometric phases that modify the angular
coherence.
The detected field is a weighted superposition over rays labelled by direction
offset $\boldsymbol{\theta}$:
\begin{equation}
  E_{\rm obs}(t) = \int d\boldsymbol{\theta}\,I^{1/2}(\boldsymbol{\theta})\,
  E_{\boldsymbol{\theta}}\!\left(t - \delta t(t;\boldsymbol{\theta})\right),
  \label{eq:EobsExt}
\end{equation}
where the $E_{\boldsymbol{\theta}}$ are independent, identically distributed
complex Gaussian processes (one per sky direction), and each ray carries its
own foam jitter $\delta t(t;\boldsymbol{\theta})$.
The spatiotemporal covariance of the path-length aggregate is given by
Eq.~\eqref{eq:pathcov} of Section~\ref{sec:model}, with
$b(\boldsymbol{\theta},\boldsymbol{\theta}') \simeq D|\boldsymbol{\theta}-\boldsymbol{\theta}'|$
the transverse separation between rays evaluated at the source end.
For foam distributed uniformly along the line of sight, the ray separation varies as $z|\boldsymbol{\theta}-\boldsymbol{\theta}'|$; evaluating at the source plane $b = D|\boldsymbol{\theta}-\boldsymbol{\theta}'|$ introduces an $O(1)$ geometric factor that is absorbed into the free parameter $r_\perp$.

The first-order coherence is unchanged from the point-source case:
spatial incoherence between source directions
($\langle E^*_{\boldsymbol{\theta}}\,E_{\boldsymbol{\theta}'}\rangle
\propto \delta(\boldsymbol{\theta}-\boldsymbol{\theta}')$) forces only the
diagonal to contribute, and the angular integral over $I$ gives unity, so
\begin{equation}
  g^{(1)}_{\rm obs}(\tau)
  = g^{(1)}_{\rm src}(\tau)\cdot e^{-\sigma_\phi^2[1-C_\alpha(\tau;D)]},
  \label{eq:g1ext}
\end{equation}
identical to Eq.~\eqref{eq:g1narrowband}. The transverse foam structure
leaves no imprint on $g^{(1)}$.

It is convenient to define the effective transverse coherence factor
weighting $\Phi(b/r_\perp)$ by the source brightness distribution:
\begin{equation}
  \Phi_{\rm eff}
  \equiv \int\!\!\int d\boldsymbol{\theta}\,d\boldsymbol{\theta}'\;
         I(\boldsymbol{\theta})I(\boldsymbol{\theta}')\,
         \Phi\!\left(\frac{b(\boldsymbol{\theta},\boldsymbol{\theta}')}{r_\perp}\right),
  \label{eq:PhiEff}
\end{equation}
with limits $\Phi_{\rm eff}\to 1$ (point-like source, or
$r_\perp\gg D\,\Delta\theta$) and $\Phi_{\rm eff}\to 0$ (fully resolved,
$r_\perp\ll D\,\Delta\theta$).

\subsection{Second-Order Coherence: Connected Wick Pairing}
\label{sec:g2ext}

The unnormalized intensity correlator $\langle I(t)\,I(t+\tau)\rangle$ is a
four-point function in $E_{\rm obs}$.  Conditioning on the realization of
$\delta t$ and applying Wick's theorem to the independent zero-mean
complex-Gaussian $E_{\boldsymbol{\theta}}$ produces two pairings; spatial
incoherence then collapses the four-fold angular integral to a double
integral over ray pairs.  The disconnected pairing
($\boldsymbol{\theta}_1=\boldsymbol{\theta}_2$,
$\boldsymbol{\theta}_3=\boldsymbol{\theta}_4$) contributes
$\langle I(t)\rangle\langle I(t+\tau)\rangle$ and cancels in $g^{(2)}-1$.
The connected pairing forces $\boldsymbol{\theta}_1=\boldsymbol{\theta}_4
\equiv\boldsymbol{\theta}$ and $\boldsymbol{\theta}_3=\boldsymbol{\theta}_2
\equiv\boldsymbol{\theta}'$, and yields two surviving two-point functions
evaluated at retarded-time differences that depend on \emph{two
same-direction jitter increments}:
\begin{equation}
  X \equiv \delta t(t+\tau;\boldsymbol{\theta}) - \delta t(t;\boldsymbol{\theta}),
  \qquad
  Y \equiv \delta t(t+\tau;\boldsymbol{\theta}') - \delta t(t;\boldsymbol{\theta}').
  \label{eq:XYdef}
\end{equation}
$X$ is the temporal jitter increment along direction $\boldsymbol{\theta}$;
$Y$ is the analogous increment along direction $\boldsymbol{\theta}'$.
Neither involves a cross-direction sample of $\delta t$; the transverse
coherence enters only through their joint distribution.  Using
$g^{(1)}_{\rm src}(-u) = g^{(1)*}_{\rm src}(u)$, the connected piece is
\begin{equation}
  g^{(2)}_{\rm obs}(\tau) - 1
    = \int\!\!\int d\boldsymbol{\theta}\,d\boldsymbol{\theta}'\;
      I(\boldsymbol{\theta})I(\boldsymbol{\theta}')\;
      \big\langle g^{(1)}_{\rm src}(\tau - X)\,
                 g^{(1)*}_{\rm src}(\tau - Y)\big\rangle_{X,Y}.
  \label{eq:g2ext_correct}
\end{equation}
The full Wick contraction is carried out in Appendix~\ref{app:g2ext}.

$X$ and $Y$ are linear functionals of the Gaussian process $\delta t$,
hence jointly Gaussian and zero-mean, with
\begin{equation}
  \operatorname{Var}(X) = \operatorname{Var}(Y)
    = \frac{2\sigma_\ell^2(D)}{c^2}\big[1 - C_\alpha(\tau;D)\big], \qquad
  \operatorname{Cov}(X,Y)
    = \frac{2\sigma_\ell^2(D)}{c^2}\big[1 - C_\alpha(\tau;D)\big]\,
       \Phi(b/r_\perp),
  \label{eq:XYvar}
\end{equation}
so the correlation coefficient is $\rho = \Phi(b/r_\perp)$.

Equation~\eqref{eq:g2ext_correct} is a 2D average over the joint
distribution of $(X,Y)$, not a 1D convolution with a single jitter
variable.  It collapses to a 1D convolution only at $\Phi = 1$ (point
source, or $b\ll r_\perp$): there $X = Y$ almost surely,
$g^{(1)}_{\rm src}(\tau-X)\,g^{(1)*}_{\rm src}(\tau-X) =
|g^{(1)}_{\rm src}(\tau-X)|^2 = g^{(2)}_{\rm src}(\tau-X) - 1$, and the
point-source formula Eq.~\eqref{eq:g2obs} is recovered with
$\sigma_\Delta^2 = \operatorname{Var}(X) =
(2\sigma_\ell^2/c^2)[1-C_\alpha(\tau)]$.

\subsection{Finite-Lag Suppression}
\label{sec:zerolag}

\paragraph{Zero lag.}
At $\tau=0$ each of the increments in Eq.~\eqref{eq:XYdef} is the
difference of a random variable with itself: $X = 0$ and $Y = 0$
\emph{identically}, for every realization of $\delta t$ and every
$(\boldsymbol{\theta},\boldsymbol{\theta}')$ --- independent of $\Phi$,
$C_\alpha$, $\sigma_\ell$, and $D$.  Equation~\eqref{eq:g2ext_correct}
reduces to $\int\!\!\int I\,I'\,|g^{(1)}_{\rm src}(0)|^2 = 1$, so
\begin{equation}
  g^{(2)}_{\rm obs}(0) = 2 \qquad \text{exactly,}
  \label{eq:g2zero_exact}
\end{equation}
and $\mathcal{S}(0) = g^{(2)}_{\rm obs}(0) - 2 = 0$ for any extended source
in this model.  The conditional Gaussianity of $E_{\rm obs}(t)$ given
$\delta t$ (see Appendix~\ref{app:g2ext}) preserves the single-mode
bunching peak exactly; foam imprints itself on multi-time correlations
but not on the single-time amplitude statistics.

\paragraph{Finite lag.}
For $\tau \neq 0$ the joint $(X,Y)$ average develops nontrivial structure.
Splitting $g^{(1)}_{\rm src}(\tau)=e^{i\omega_0\tau}V_{\rm src}(\tau)$ into
carrier and envelope and changing variables to $U \equiv (X+Y)/2$ and
$V \equiv X-Y$ (independent Gaussians of variances
$\sigma_U^2 = (\sigma_\ell^2/c^2)[1-C_\alpha](1+\Phi)$ and
$\sigma_V^2 = (4\sigma_\ell^2/c^2)[1-C_\alpha](1-\Phi)$; see
Appendix~\ref{app:g2ext}), the narrow-band reduction yields
\begin{equation}
  g^{(2)}_{\rm obs}(\tau) - 1
  = \int\!\!\int d\boldsymbol{\theta}\,d\boldsymbol{\theta}'\,
    I(\boldsymbol{\theta})I(\boldsymbol{\theta}')\,
    \mathcal{F}_{\rm cross}(\tau;b)\,
    \int dU\,\mathcal{N}(U;0,\sigma_U^2)\,V_{\rm src}^2(\tau-U),
  \label{eq:g2ext_corrected}
\end{equation}
with the \emph{cross-direction carrier-phase factor}
\begin{equation}
  \mathcal{F}_{\rm cross}(\tau;b)
  \equiv \exp\!\Big(-2\sigma_\phi^2(D)\,[1-C_\alpha(\tau;D)]\,
                    [1-\Phi(b/r_\perp)]\Big),
  \label{eq:Fcross}
\end{equation}
which reduces to unity at $\tau=0$ (recovering Eq.~\eqref{eq:g2zero_exact}),
at $\Phi=1$ (recovering the point-source formula
Eq.~\eqref{eq:g2obs}), and in the quasi-static foam limit $C_\alpha\to 1$.
The decoupling between the imaging-channel scale $\sigma_\phi$ and the
HBT scale $\sigma_\Delta/\tau_c$ established in
Section~\ref{sec:scales} is exact only for a point source: an extended
source viewed at finite lag picks up $\sigma_\phi$ in $g^{(2)}$ through
$\mathcal{F}_{\rm cross}$.

\paragraph{Siegert discrepancy.}
Inserting Eqs.~\eqref{eq:g2ext_corrected} and \eqref{eq:g1ext} into the
Siegert-discrepancy definition $\mathcal{S}(\tau) \equiv [g^{(2)}_{\rm obs}(\tau)-1]
- |g^{(1)}_{\rm obs}(\tau)|^2$, expanding in the weak-foam regime
$\sigma_\phi^2[1-C_\alpha(\tau)] \ll 1$, and performing the angular
integral gives
\begin{equation}
  \mathcal{S}_{\rm ext}(\tau)
  \approx \Phi_{\rm eff}\cdot\mathcal{S}_{\rm pt}(\tau),
  \label{eq:SiegertExt}
\end{equation}
where $\mathcal{S}_{\rm pt}(\tau) \approx 2\sigma_\phi^2[1-C_\alpha(\tau)]\,
V_{\rm src}^2(\tau)$ is the point-source Siegert discrepancy of
Section~\ref{sec:siegert}.  The extended-source Siegert discrepancy is
therefore \emph{suppressed} by $\Phi_{\rm eff}\in[0,1]$ relative to the
point-source value: the point-source result is the upper bound on the
foam signature in this framework.  Targets believed to be intrinsically
compact ($D\Delta\theta \ll r_\perp$ for any plausible $r_\perp$) saturate
the upper bound and are the safest detection candidates; partially
resolved targets degrade the SNR by the same factor $\Phi_{\rm eff}$,
and the angular-size dependence is itself a positive discriminant ---
comparing $\mathcal{S}(\tau)$ across targets of varying angular size at
fixed $D$ probes $r_\perp$ through the profile $\Phi(b/r_\perp)$.

\subsection{Discriminants Against Conventional Phase Screens}

The Siegert observable serves as a discriminant against conventional
propagation phase screens~\cite{Carlip2023}; for foam models with
long-range transverse correlations --- where image blurring is strongly
suppressed~\cite{LeeZurekChen2024} --- it may be the only sensitive
channel.  The mathematical structure of the foam suppression factor
$\mathcal{F}_{\rm foam}$ is identical to that of any Gaussian random
phase screen, but two signatures distinguish foam: distance scaling and
achromaticity.  For spacetime foam
$\sigma_\phi^2\propto D^{2(1-\alpha)}$ ($D^1$ for $\alpha=1/2$,
$D^{2/3}$ for $\alpha=2/3$, $D^0$ for $\alpha=1$), whereas a thin
intervening screen scales as $D$ and Kolmogorov turbulence as
$D^{11/6}$; the holographic and conservative foam models predict
sub-linear or flat growth qualitatively different from any known
astrophysical screen.  Geometrically, foam jitter
($\delta t=\Delta\ell/c$) is achromatic --- independent of $\lambda$ ---
in contrast to plasma-based effects which scale as $\lambda^2$.  The
Fermi, H.E.S.S., and LHAASO time-of-flight analyses~\cite{Fermi2009,HESS2008,LHAASO2024}
rule out deterministic, energy-ordered LIV dispersion as a class of
quantum-gravity effect but do not constrain the stochastic achromatic
foam jitter considered here (Section~\ref{sec:intro}); the principal
chromatic alternative the achromaticity discriminant addresses is
therefore plasma propagation.  In practice, the intrinsic source
visibility $V_{\rm src}(\tau)$ introduces a degeneracy with foam
suppression, but the Siegert discrepancy $\mathcal{S}(\tau)$ is
constructed directly from the data: a detection of $\mathcal{S}\neq 0$
requires no explicit model of the source spectrum, since
$g^{(1)}_{\rm src}(\tau)$ can be obtained from the source spectrum via
the Wiener--Khinchin theorem.

% ==============================================================
\section{Siegert-Violation: Sensitivity and Target Selection}
\label{sec:prospects}
% ==============================================================

This section examines the detectability of Siegert violation via the HBT
split-beam configuration~\cite{HBT1956,Foellmi2009,Dravins2013,Guerin2017}
and surveys candidate sources with Lorentzian (or Lorentzian-dominated)
emission lines for which the sensitivity estimates apply directly. The
single foam observable is the finite-lag Siegert discrepancy
$\mathcal{S}(\tau)$ of Section~\ref{sec:siegert}; the point-source case
of Section~\ref{sec:sens-point} gives its maximum-signal limit, and the
extended-source case of Section~\ref{sec:sens-extended} reports the same
observable scaled by $\Phi_{\rm eff}$.  We treat the resulting SNR
figures as a statistical floor; in practice systematic errors will set
the binding limit and a dedicated systematic budget is required for any
claimed detection.

\subsection{Sensitivity and Aperture Scaling}

For a single-aperture HBT configuration using a 50/50 beamsplitter with
unpolarized light~\cite{Guerin2025}, the per-lag noise on $g^{(2)}_{\rm obs}$ is
\begin{equation}
  \sigma[g^{(2)}_{\rm obs}]
  \;=\; \frac{4}{R\sqrt{\tau_r\,T_{\rm obs}}}\,,
  \label{eq:g2noise}
\end{equation}
where $R$ is the total incident photon count rate, $\tau_r$ the timing resolution
(correlation bin width) of the detectors, and $T_{\rm obs}$ the total observation
time.
The prefactor $4 = 2N_p$ with $N_p=2$ polarization modes reflects that only
photons sharing a polarization state contribute to the bunching signal, while
the shot noise is set by the total (polarization-summed) count rate.
The SNR for detecting a Siegert discrepancy of amplitude $\mathcal{S}(\tau)$
evaluated at a specific lag $\tau$ --- treating $\mathcal{S}$ as constant
within the bin of width $\tau_r$ --- is then
\begin{equation}
  \mathrm{SNR}_{\mathcal{S}}
  \;=\; \frac{\mathcal{S}}{\sigma[g^{(2)}_{\rm obs}]}
  \;=\; \frac{\mathcal{S}\,R}{4}\sqrt{\tau_r\,T_{\rm obs}}\,.
  \label{eq:snr_siegert}
\end{equation}

All numerical projections adopt $A_\alpha = 1$ and the dynamical-foam limit
$C_\alpha(\tau_L; D) \to 0$.  The Siegert observable requires
$\tau_{\rm foam}\lesssim\tau_L$ so that $1-C_\alpha(\tau_L)$ is of order
unity; in the quasi-static limit $C_\alpha\to 1$ the signal vanishes for
all targets and the channel becomes inoperative.

The photon-counting rate is $R = \epsilon\, A\, \dot{n}$, where $\dot{n}$ is the
source photon flux per unit area within the filter bandpass, $\epsilon$ the
end-to-end throughput, and $A$ the collecting area. We adopt the Kim
\emph{et al.}~\cite{Kim2025} baseline: $d = 9.96$\,m
($A = 77.9$\,m$^2$), $\epsilon = 0.39$, and $\sigma_t = 13$\,ps rms
(30\,ps FWHM).
The astrophysical targets of interest have intrinsic coherence times
$\tau_L \lesssim 1$\,ps, so $\sigma_t \gg \tau_L$ for the baseline instrument and
the measurement sits firmly in the detector-limited regime: an unresolved
foam feature of width $\tau_L$ in a bin of width $\sigma_t$ contributes an
effective signal $(\tau_L/\sigma_t)\,|\mathcal{S}|$, so
Eq.~\eqref{eq:snr_siegert} becomes
\begin{equation}
  \mathrm{SNR}_{\rm det}
  \;=\; \frac{|\mathcal{S}|\,R\,\tau_L}{4}\sqrt{\frac{T_{\rm obs}}{\sigma_t}}\,.
  \label{eq:snr_det}
\end{equation}
We present the detector-limited SNR as the primary case throughout and quote
the source-limited limit ($\tau_L \gtrsim \sigma_t$, requiring sub-ps timing)
only as the gain that next-generation detectors would unlock.

Achieving a statistically significant SNR is a necessary but not sufficient
condition for a foam detection. A separate and equally demanding challenge
is controlling systematic errors---spurious correlations from detector
afterpulsing, electronic crosstalk, and optical back-reflections produce
positive artifacts near zero lag, atmospheric scintillation introduces
classical intensity fluctuations that mimic $g^{(2)}>1$, and detector
dead time causes a compensating dip in $g^{(2)}$ near zero lag---to levels
below the expected foam signal $|\mathcal{S}|$. Because the foam signatures
are small ($|\mathcal{S}|\ll 1$ for all viable foam models), even modest
systematics can mimic or mask the effect.

\subsection{Finite-Lag Channel: Point Sources}
\label{sec:sens-point}

For a point-like Lorentzian source in the fast-foam / dynamical limit
($\tau_{\rm foam}\ll\tau_L$, $C_\alpha\to 0$ above $\tau_{\rm foam}$), the
Siegert discrepancy evaluated at its peak lag $\tau_{\max}$
(Eq.~\eqref{eq:taumax_exp}) is given by Eq.~\eqref{eq:Stc},
$\mathcal{S}(\tau_{\max}) \approx 1 - e^{-2\sigma_\phi^2}$ with
$\sigma_\phi$ the rms accumulated phase (Eq.~\eqref{eq:sigmaphi}): quadratic
in $\sigma_\phi$ in the weak-foam limit
($\mathcal{S}(\tau_{\max})\approx 2\sigma_\phi^2$) and saturating at
$1$ for $\sigma_\phi\gg 1$.

\paragraph{Detector-limited SNR.}
Substituting $\mathcal{S}(\tau_{\max})$ into Eq.~\eqref{eq:snr_det},
\begin{equation}
\textrm{SNR}_{\rm det}^{\rm pt}
\;=\; \frac{|\mathcal{S}(\tau_{\max})|\,R\,\tau_L}{4}\sqrt{\frac{T_{\rm obs}}{\sigma_t}}
\;\xrightarrow{\sigma_\phi\ll 1}\;
\tfrac{1}{2}\sigma_\phi^2\, R\,\frac{\tau_L}{\sqrt{\sigma_t}}\sqrt{T_{\rm obs}}.
\label{eq:snr_det_pt}
\end{equation}
The $\tau_L$ factor here is the source coherence width (the natural HBT
coincidence-bin scale), not the evaluation lag.
The detector-limited form carries a $\sqrt{\tau_L/\sigma_t}$ penalty relative
to source-limited operation
($\textrm{SNR}_{\rm src}^{\rm pt} = \tfrac{1}{4}\,|\mathcal{S}(\tau_{\max})|\,R\sqrt{\tau_L T_{\rm obs}}$);
reaching the source-limited limit on broad-line quasars requires
$\sigma_t \lesssim \tau_L \sim 0.1$\,ps.

\paragraph{Filter strategy.}
The foam signal $\mathcal{S}(\tau_{\max})$ depends only on $\sigma_\phi$, not on
the filter width.
The SNR (Eq.~\eqref{eq:snr_det_pt}) scales as $R\,\tau_L$, and these two
quantities trade off against each other as the filter is narrowed: a narrower
filter reduces the transmitted flux $R$ but increases the post-filter
coherence time $\tau_c = \lambda^2/(c\,\Delta\lambda_{\rm filter})$.
For a Lorentzian emission line, $R\propto\Delta\lambda_{\rm filter}$ and
$\tau_c\propto\Delta\lambda_{\rm filter}^{-1}$ in the narrow-filter limit,
so $R\,\tau_c$ is constant and the SNR is insensitive to filter width.
The same cancellation holds for a flat continuum source.
The practical lower bound on filter width is sky background: widening the
filter beyond the emission line adds sky photons without increasing the
source signal.

\paragraph{Distance scaling.}
Substituting $R\propto D^{-2}$ at fixed luminosity and
$\sigma_\phi^2\propto D^{2(1-\alpha)}$ into Eq.~\eqref{eq:snr_det_pt} gives, in
the weak-foam limit,
\begin{equation}
\textrm{SNR}^{\rm pt} \;\propto\; D^{2(1-\alpha)}\cdot D^{-2} \;=\; D^{-2\alpha},
\label{eq:dist_scaling_pt}
\end{equation}
saturating at $\textrm{SNR}\propto D^{-2}$ in the strong-foam limit.
Since $\alpha > 0$, the SNR increases toward smaller $D$: nearby targets are
preferred, provided the flat-space ($z\ll 1$) limit of
Eq.~\eqref{eq:sigmaphi} holds.
At cosmological distances the $D$-dependence is modified by two competing
effects.
For a fixed rest-frame line, $\lambda_{\rm obs}=\lambda_{\rm rest}(1+z)$ grows
with redshift, reducing $\sigma_\phi$ via the $1/\lambda_{\rm obs}^2$ prefactor
in Eq.~\eqref{eq:sigmaphi}.
Conversely, a fixed observer-frame bandpass accesses progressively bluer
rest-frame transitions at higher redshift
($\lambda_{\rm rest}=\lambda_{\rm obs}/(1+z)$); since
$\sigma_\phi\propto 1/\lambda_{\rm rest}$, a source whose spectral line falls
in the UV rest-frame accumulates substantially more phase than a lower-redshift
counterpart observed at the same detector wavelength---potentially reversing the
$D^{-2\alpha}$ preference.

The minimum detectable path-length amplitude follows from inverting
Eq.~\eqref{eq:snr_det_pt} with the quadratic signal dependence
$|\mathcal{S}(\tau_{\max})|\propto\sigma_\phi^2\propto\sigma_\ell^2$:
\begin{equation}
\sigma_{\ell,\rm min}^{\rm pt} \;\propto\;
\frac{\lambda}{d}\!\left(\frac{\sigma_t}{\tau_L^2\, T_{\rm obs}}\right)^{\!1/4},
\label{eq:sigma_ell_min_pt}
\end{equation}
where $R\propto d^2$ has been used. Sensitivity improves as $d^{-1}$ rather than
$d^{-2}$, and as $T_{\rm obs}^{-1/4}$ rather than
$T_{\rm obs}^{-1/2}$ --- both costs of the quadratic dependence on
$\sigma_\phi$. Each additional decade of sensitivity in $\sigma_\ell$ requires
four decades of integration time, compared with two for the extended-source
variant.

\subsection{Finite-Lag Channel: Extended Sources}
\label{sec:sens-extended}

For an extended source viewed at finite lag, the Siegert discrepancy is
suppressed relative to the point-source value by the angular-size factor
$\Phi_{\rm eff}$ of Eq.~\eqref{eq:PhiEff}.  In the weak-foam regime
relevant to all viable foam models, Eq.~\eqref{eq:SiegertExt} gives
\begin{equation}
  \mathcal{S}_{\rm ext}(\tau_{\max})
  \;\approx\; \Phi_{\rm eff}\cdot\mathcal{S}_{\rm pt}(\tau_{\max})
  \;\approx\; 2\,\Phi_{\rm eff}\,\sigma_\phi^2,
  \label{eq:SiegertExtPeak}
\end{equation}
and the detector-limited SNR follows by scaling Eq.~\eqref{eq:snr_det_pt}:
\begin{equation}
  \textrm{SNR}_{\rm det}^{\rm ext}
  \;=\; \Phi_{\rm eff}\cdot\textrm{SNR}_{\rm det}^{\rm pt}.
  \label{eq:snr_det_ext}
\end{equation}
The point-source SNR of Section~\ref{sec:sens-point} is the
\emph{maximum-signal upper bound} for the foam observable in this
framework; no target gives SNR above the point-source value.  The
extended-source factor $\Phi_{\rm eff}\in[0,1]$ scales this upper bound
down for partially resolved targets and depends on the unknown foam
transverse coherence length $r_\perp$:
$\Phi_{\rm eff}\to 1$ when $D\Delta\theta \ll r_\perp$ (the source is
unresolved by the foam, recovering the upper bound), and
$\Phi_{\rm eff}\to 0$ when $D\Delta\theta \gg r_\perp$ (the source is
fully resolved, signal vanishes).

\paragraph{Caveat on $r_\perp$.}
Because $r_\perp$ is an unknown foam-model parameter, the
extended-source SNR is reported as a band over plausible $r_\perp$
rather than as a single number.  Targets believed to be intrinsically
compact (e.g., narrow-line AGN cores with $D\Delta\theta$ smaller than
any plausible $r_\perp$) saturate the point-source bound and provide
the most conservative detection candidates; partially resolved targets
degrade the SNR by $\Phi_{\rm eff}$, and a non-detection in such a
target is consistent either with foam below threshold or with $r_\perp$
small enough to resolve the source.

\paragraph{Angular-size discriminant.}
Should a signal be detected, the source-size dependence of
$\mathcal{S}(\tau)$ across targets of varying $D\Delta\theta$ is a
positive observational handle on $r_\perp$, with no analogue in a
featureless propagation screen (Section~\ref{sec:extended}).  A
multi-target campaign probes $r_\perp$ through the profile
$\Phi(b/r_\perp)$ even when no individual measurement does.

\paragraph{Distance scaling.}
The distance scaling of Eq.~\eqref{eq:dist_scaling_pt} is unchanged for
extended targets up to the $\Phi_{\rm eff}$ prefactor:
$\textrm{SNR}^{\rm ext}\propto \Phi_{\rm eff}\cdot D^{-2\alpha}$ in the
weak-foam limit.  Filter strategy, $\tau_L$ scaling, and the minimum
detectable path-length amplitude all carry over directly from
Section~\ref{sec:sens-point} with an overall $\Phi_{\rm eff}$ factor in
the signal.

\subsection{Candidate Sources}
\label{sec:candidate_sources}

Two source properties drive target selection: the photon rate $R$, set by
the total line flux captured within the filter, and the source distance
$D$, which sets $\sigma_\phi$ via Eq.~\eqref{eq:sigmaphi} and hence the
signal amplitude $\mathcal{S}(\tau_{\max}) \approx 1-e^{-2\sigma_\phi^2}$
of Eq.~\eqref{eq:Stc}.  For partially resolved targets the signal is
further suppressed by $\Phi_{\rm eff}$ (Section~\ref{sec:sens-extended}),
which depends on the unknown foam transverse coherence length $r_\perp$;
the entries below quote the point-source upper bound and flag $\Phi_{\rm eff}$
caveats where the source has substantial angular extent.
All sources discussed below have intrinsic coherence times $\tau_L \lesssim 1$\,ps,
well below the baseline timing resolution $\sigma_t = 13$\,ps; the measurements
therefore operate in the detector-limited regime throughout.
The source classes surveyed below are illustrative rather than exhaustive;
systematic searches through existing AGN flux catalogues and emission-line
surveys will likely identify superior targets with higher photon rates,
more purely Lorentzian line profiles, or more favorable distance-to-flux ratios.
Observed properties and references for the principal candidates are collected
in Table~\ref{tab:source_agn}.

\begin{table*}[t]
\centering
\caption{Candidate sources for the sensitivity analysis of
Section~\ref{sec:prospects}.
$D$ is the luminosity distance (flat $\Lambda$CDM, Planck
2018~\cite{Aghanim2020}); $m$ is the apparent magnitude (band per
source class); $F$ is the integrated line flux; FWHM is observed-frame;
$\tau_L = \lambda_{\rm obs}^2/(\pi c\,\Delta\lambda_{\rm FWHM})$.
Per-source caveats (sample-mean SN~IIn scaling, $\eta$\,Car
factor-of-a-few uncertainty, J1330$-$0905 image-A flux derivation) are
given in Section~\ref{sec:candidate_sources}.
$^{\dagger}$Adopted from the Kollatschny~\&~Zetzl~2013 turbulent
FWHM~\cite{Kollatschny2013}; not resolved by the XP spectrum.}
\label{tab:source_agn}
\begin{tabular}{llcccccc}
\toprule
Source & $z$ & $D$ (Mpc) & $m$ & Line & $F$ (erg\,s$^{-1}$\,cm$^{-2}$) & FWHM (\AA) & $\tau_L$ (ps) \\
\midrule
Mrk~766          & 0.013 &  59                 & 13.6 & H$\alpha$ & $2.8\times10^{-12}$~\cite{RodriguezArdila2006}                    &  28~\cite{Kollatschny2013} & 0.17  \\
PHL\,1811        & 0.192 & 971                 & 14.4 & H$\alpha$ & $1.0\times10^{-12}$~\cite{Leighly2007}                            &  46~\cite{Leighly2007}     & 0.14  \\
$\eta$\,Car (1843) & 0   & $2.3\times10^{-3}$~\cite{2001ApJ...553..837H}  & $-1$   & H$\alpha$ & ${\sim}5\times10^{-7}$~\cite{2018MNRAS.480.1466S}                  &  15~\cite{2018MNRAS.480.1466S} & 0.30  \\
Gal.\ SN~IIn     & 0     & $2\times10^{-3}$    & $-7.6$ & H$\alpha$ & ${\sim}2\times10^{-4}$~\cite{Taddia2013,2025ApJ...987...13R} &  55~\cite{Taddia2013} & 0.083 \\
M31 SN~IIn       & 0     & $0.78$              & $5.3$  & H$\alpha$ & ${\sim}1.4\times10^{-9}$~\cite{Taddia2013,2025ApJ...987...13R} &  55~\cite{Taddia2013} & 0.083 \\
J1330$-$0905 (A) & 2.225 & $1.81\times10^{4}$  & $16.1$ & Ly$\alpha$ & $5.1\times10^{-13}$~\cite{PersephonesTorch} &  52$^{\dagger}$        & 0.031 \\
\bottomrule
\end{tabular}
\end{table*}

\textit{Narrow-line Seyfert 1 galaxies as headline AGN targets.}
NLS1s (Population~A; FWHM $\lesssim 4000$\,km\,s$^{-1}$;
\cite{2002ApJ...566L..71S,2010MNRAS.403.1759Z}) have single-Lorentzian
broad lines and intrinsically compact broad-line regions, both
favourable for the Siegert observable.  The NLS1 BLR has angular size
$\sim 10^{-6}$~arcsec at cosmological distances, far below any
plausible $r_\perp$, so $\Phi_{\rm eff}\approx 1$ and these targets
saturate the point-source upper bound.
We adopt two reference targets at complementary distances --- the nearby
Mrk~766 and the cosmological PHL~1811 (Table~\ref{tab:source_agn};
luminosity distances in flat $\Lambda$CDM with the Planck~2018
cosmology~\cite{Aghanim2020}); the H$\alpha$ and H$\beta$ profiles of
PHL~1811 are explicitly Lorentzian-fit by Leighly~et~al.~2007 (their
\S3)~\cite{Leighly2007}.
A wide filter ($\Delta\lambda_{\rm filter} \geq \mathrm{FWHM}$) captures
the full line flux, giving photon rates $R \approx 2.0\times 10^5$\,cps
(Mrk~766) and $R \approx 8.5\times 10^4$\,cps (PHL~1811) on a 10\,m
telescope with $\epsilon = 0.39$; sky background and detector dark
counts are negligible relative to the source rates, so the Brouri~et~al.\
correction for accidental coincidences is small~\cite{Brouri2000}.
For both targets the $2\times \mathrm{FWHM}_{\rm line}$ filter width is
narrow enough that the underlying pseudo-continuum (Fe~II forest plus
near-line continuum) dilutes the signal by $\lesssim 10\%$ given
the line-to-continuum ratio $\gtrsim 10$:1 typical of these objects;
the dilution does not change the qualitative SNR ordering but should be
subtracted from the measured $g^{(2)}(\tau)$ as an additive constant
when extracting $\mathcal{S}(\tau)$ in practice.
The two targets span a factor of 16 in distance, exercising the
$\sigma_\phi \propto D^{1-\alpha}$ scaling of the foam-vs-screen
discriminant developed in Section~\ref{sec:extended}.

\textit{Strongly-lensed quasars.}
Gravitational magnification by factors of ${\sim}10$--$60$ boosts the photon
rate while preserving the cosmological foam-variance lever arm at the
source's intrinsic distance, making strongly-lensed quasars a complementary
persistent-AGN target class.
Persephone's Torch (J1330$-$0905, $z=2.22$, total magnification
${\sim}56$, brightest-image $i\approx 16$)~\cite{PersephonesTorch} is
included in Table~\ref{tab:source_agn} on the basis of its broad
Ly$\alpha$ emission line measured directly from its GAIA~DR3 XP spectrum
(observed equivalent width 235\,\AA, integrated 3700--4100\,\AA\ line
flux $1.28\times10^{-12}$~erg\,s$^{-1}$\,cm$^{-2}$ for the GAIA blob;
image~A carries a fraction 0.40 of the total system flux from the
observed ratios $F_B/F_A = 0.74$, $F_C/F_A = 0.70$, $F_D/F_A = 0.08$,
giving $5.1\times10^{-13}$~erg\,s$^{-1}$\,cm$^{-2}$ on image~A).
Under the Kollatschny~\&~Zetzl~2013~\cite{Kollatschny2013}
broad-line-region model the observed Ly$\alpha$ profile is a Voigt
profile with a Lorentzian core from BLR turbulent broadening
(FWHM$_{\rm turb} \sim 3800$\,km\,s$^{-1}$) convolved with a rotational
Gaussian; the Lorentzian core sets the coherence time $\tau_L$ relevant
for the finite-lag Siegert observable.  The BLR angular size at $z=2.2$
is $\sim 10^{-7}$~arcsec, far below any plausible $r_\perp$, so
$\Phi_{\rm eff}\approx 1$ and this target saturates the point-source
upper bound.
The net finite-lag SNR exceeds PHL\,1811 by ${\sim}1.4\times$ at
$\alpha = 2/3$ (Figure~\ref{fig:snr_vs_alpha_pt}).
The Lorentzian-core hypothesis at high $z$ is an extrapolation from the
low-$z$ K\&Z sample; direct high-resolution spectroscopy of Ly$\alpha$ (or
the redshifted near-infrared Balmer lines) would confirm it.

\textit{Galactic and Local-Group transients.}
Type~IIn supernovae and LBV giant eruptions --- with $\eta$\,Car as the
principal example --- combine high H$\alpha$ luminosity with narrow
Lorentzian line cores ($\tau_L \sim 0.1$--$0.3$\,ps).  The point-source
upper-bound SNR scales as $\sigma_\phi^2\propto D^{2(1-\alpha)}$ and is
small at kpc--Mpc distances; however, the giant photon rate $R$ at close
range partially compensates, and the resulting SNR remains competitive
for $\alpha\to 1$ (conservative foam).  At smaller $\alpha$ these
nearby targets lose the cosmological lever arm and fall below the AGN
baselines (Section~\ref{sec:snr_results}).
Two additional target-specific caveats apply.
\textit{Source extent:} a Galactic SN~IIn or LBV ejecta photosphere has
angular extent $\sim 10^{-5}$--$10^{-3}$~arcsec, which can exceed
$D\Delta\theta\sim r_\perp$ for small $r_\perp$; the SNR is then
suppressed by $\Phi_{\rm eff}(r_\perp)$ and the point-source bound is
not necessarily saturated.
\textit{Scattering geometry:} the Homunculus Nebula around $\eta$\,Car
scatters substantial flux into the line of sight; the observed coherence
is partially set by scattering geometry and must be modeled before
crediting the photon count.
\textit{Stationarity:} the framework assumes a stationary Gaussian source,
valid for integrations short relative to the variability timescale
(weeks-to-months for LBV eruptions and SN~IIn light-curves); the
$10$\,hr per-event integrations used in the figures satisfy this
comfortably, while $10^4$\,hr campaign-scale totals require co-adding
over stationary epochs (Section~\ref{sec:snr_results}).

\textit{Type IIn supernovae.}
Rather than a specific event, we adopt a representative Type~IIn source
characterized by the sample-mean peak H$\alpha$ luminosity
$L({\rm H}\alpha) = 10^{41}$\,erg\,s$^{-1}$~\cite{Taddia2013}, broad-line
FWHM $\approx 55$\,\AA\ (H$\alpha$ rest-frame)~\cite{Taddia2013}, and
apparent $r$-band magnitude from the sample-mean absolute magnitude
$M_r = -19.16$~\cite{2025ApJ...987...13R}.
The H$\alpha$ flux at a given distance follows by inverse-square scaling;
we evaluate two cases: a Galactic event at $D = 2$\,kpc and an M31 event
at $D = 778$\,kpc (Table~\ref{tab:source_agn}).

\begin{figure}[t]
  \centering
  \includegraphics[width=0.92\linewidth]{snr_vs_alpha_pointsource.png}
  \caption{%
    Maximum-signal (point-source) finite-lag Siegert discrepancy SNR for
    $\mathcal{S}(\tau_{\max}) \approx 1 - e^{-2\sigma_\phi^2}$
    (Eq.~\eqref{eq:Stc}) as a function of foam exponent $\alpha$, in the
    fast-foam / dynamical limit $C_\alpha \to 0$.
    The values shown are the upper bound for each target; partially
    resolved targets ($D\Delta\theta\gtrsim r_\perp$) are suppressed by
    $\Phi_{\rm eff}\in[0,1]$ (Section~\ref{sec:sens-extended}).
    Kim~et~al.~\cite{Kim2025} baseline (10\,m, $\epsilon = 0.39$,
    $\sigma_t = 13$\,ps rms), $T_{\rm obs} = 10$\,hr per exposure (multiply
    SNR by $\sqrt{10^3} \approx 31.6$ for the $10^4$\,hr campaign quoted
    in Sections~\ref{sec:snr_results} and~\ref{sec:conclusions});
    detector-limited regime (Eq.~\eqref{eq:snr_det_pt}).
    Filled circles mark the $\sigma_\phi = 0.3$ weak-foam validity
    boundary; curves to the right saturate at $1$.
    Vertical dashed lines: $\alpha = 0.54$ (Lieu--Hillman imaging
    exclusion at $D = 1$\,Gpc, $\lambda = 500$\,nm) and
    $\alpha = 2/3$ (holographic).
    J1330$-$0905 is evaluated on the GAIA~DR3 XP-measured Ly$\alpha$
    flux with the Kollatschny~\&~Zetzl~2013 turbulent-Lorentzian core
    (Section~\ref{sec:candidate_sources}).
    Dashed curves (right axis): $|\mathcal{S}(\tau_{\max})|$.%
    \label{fig:snr_vs_alpha_pt}}
\end{figure}

\subsection{SNR Results}
\label{sec:snr_results}

The maximum-signal (point-source) finite-lag Siegert discrepancy SNR as
a function of the foam exponent $\alpha$ is shown in
Fig.~\ref{fig:snr_vs_alpha_pt} for the six sources of
Table~\ref{tab:source_agn} at the Kim~et~al.~\cite{Kim2025} baseline
($d=9.96$\,m, $\epsilon=0.39$, $\sigma_t=13$\,ps rms) and
$T_{\rm obs}=10$\,hr per exposure (multiply by $\sqrt{10^3}\approx 31.6$
for a $10^4$\,hr campaign).  For partially resolved targets the SNR is
suppressed by the angular-size factor $\Phi_{\rm eff}\in[0,1]$, which
depends on the unknown foam transverse coherence length $r_\perp$
(Section~\ref{sec:sens-extended}); the values plotted are upper bounds
saturated by intrinsically compact targets (NLS1 cores, lensed-quasar
BLRs) and degraded by $\Phi_{\rm eff}$ for nearby transients with
substantial angular extent.

\paragraph{Strong-foam regime ($\alpha\lesssim 0.54$).}
At and below the Lieu--Hillman imaging exclusion the signal saturates at
$\mathcal{S}(\tau_{\max})\to 1$ (Eq.~\eqref{eq:Stc}), producing the
flattening of every SNR curve visible in Fig.~\ref{fig:snr_vs_alpha_pt}
at small $\alpha$. On the plateau the SNR is set entirely by the photon
rate $R$ and is independent of $\alpha$. For the highest-$R$ targets
the saturation level itself is detectable in $T_{\rm obs}=10$\,hr: a
Galactic SN~IIn at $\alpha=0.54$ reaches SNR${}\sim 110$, and an
$\eta$\,Car-class eruption is detectable at SNR${}\sim 1.2$ (each
modulo $\Phi_{\rm eff}$ for substantial angular extent). The HBT
framework is therefore statistically competitive with --- and
methodologically complementary to --- the imaging tests that set the
$\alpha=0.54$ boundary, with the advantage of sensitivity to foam
models with long-range transverse correlations that evade those
tests~\cite{LeeZurekChen2024}.

\paragraph{Holographic regime ($\alpha=2/3$).}
All signals and SNRs lie many orders of magnitude below unity. On
Mrk~766, $\sigma_\phi\sim 7.3\times 10^{-9}$ gives
$\mathcal{S}(\tau_{\max})\sim 1.1\times 10^{-16}$, and the SNR at
$T_{\rm obs}=10$\,hr is ${\sim}4.9\times 10^{-17}$.  The finite-lag
channel is therefore the only HBT probe in this regime and is not
competitive with the imaging upper bounds on its own; gains come from
photon-rate-leading targets (high-$R$ transients) or from
ultra-long campaigns combined with sub-ps timing.

\paragraph{Persistent AGN: Mrk~766 and PHL~1811.}
The two NLS1 targets span a factor of $16$ in distance and trade off
photon rate against cosmological lever arm. PHL~1811's $D^{1-\alpha}$
boost in $\sigma_\ell$ at $\alpha=2/3$ approximately compensates its
$D^{-2}$ flux deficit, leaving the two within ${\sim}10$\% of each
other (Mrk~766 nominally higher; see
Fig.~\ref{fig:snr_vs_alpha_pt}). Their near-degeneracy makes the pair
a useful internal consistency check on the $\sigma_\phi\propto D^{1-\alpha}$
scaling of the foam-vs-screen discriminant (Section~\ref{sec:extended}).
Both are intrinsically compact ($\Phi_{\rm eff}\approx 1$) and saturate
the upper bound.

\paragraph{Lensed quasar J1330$-$0905: largest $|\mathcal{S}|$ of the
persistent targets.}
The strongly-lensed quasar carries the cosmological lever arm to
$z=2.22$, giving $\sigma_\phi\approx 7.8\times 10^{-8}$ on image~A at
$\alpha=2/3$ --- roughly $5\times$ that of PHL~1811. The resulting
signal $|\mathcal{S}(\tau_{\max})|\sim 1.2\times 10^{-14}$ is
${\sim}25\times$ that of PHL~1811 and the largest of the six sources,
despite the lower photon rate from a single magnified image (dashed
curves, right axis of Fig.~\ref{fig:snr_vs_alpha_pt}).  As with the
nearby NLS1s, the BLR is intrinsically compact and saturates the
upper bound.  For any HBT campaign that can absorb the photon-rate
cost, J1330$-$0905 sets the high-water mark for foam signal amplitude
among persistent targets.

\paragraph{Galactic and Local-Group transients.}
For a Galactic transient at $D\sim 2$\,kpc, the modest cosmological
lever arm makes the upper-bound signal small in absolute terms, but
the giant photon rate $R$ can compensate. The Galactic SN~IIn case
yields the largest single-event upper-bound SNR
(${\sim}2\times 10^6$ times the Mrk~766 baseline at $\alpha=2/3$); the
recurrence rate makes it a ``drop-everything'' opportunity rather
than a campaign target. Smith~et~al.~2011~\cite{Smith2011} estimate
that ${\sim}9\%$ of core-collapse supernovae are Type~IIn, and
combined with the Galactic core-collapse rate of $1.6\pm 0.5$ per
century from Adams~et~al.~2013~\cite{Adams2013} this implies a Galactic
IIn rate of roughly one per 500--1500\,yr; the most recent confirmed
Galactic core-collapse supernova was Cassiopeia~A (${\sim}1680$,
Type~IIb). The $\eta$\,Car case sacrifices a factor of ${\sim}10$ in
SNR but offers a recurrence plausibly within a working career.
The M31 SN~IIn is intermediate (plotted adjacent to its Galactic
counterpart in Fig.~\ref{fig:snr_vs_alpha_pt} for direct comparison)
and degrades the Galactic SNR by the inverse-square-distance factor
$(778/2)^2\approx 1.5\times 10^5$.
All three transients have ejecta-photosphere angular sizes that may
exceed $r_\perp$ for small-$r_\perp$ foam models, in which case the
actual SNR is suppressed by $\Phi_{\rm eff}<1$ below the upper bounds
plotted; a non-detection in these targets does not by itself rule out
foam, since it is degenerate with small $r_\perp$
(Section~\ref{sec:sens-extended}).

\subsection{Sources Investigated but Not Competitive}
\label{sec:not_competitive}

The $\sigma_\phi \propto D^{1-\alpha}$ scaling means that any nearby source
suffers a phase-variance penalty no realistic photon-rate or spectrograph
multiplexing advantage can recover: at $D = 1$\,AU, $\sigma_\phi^2$ is
${\sim}10^{16}$ smaller than at $D = 1$\,Gpc for $\alpha = 1/2$ and is
utterly negligible for $\alpha = 2/3$. Three illustrative cases bear this
out.
\textit{The Sun as a continuum source.} A 1\,\AA\ filter on a 10\,m
telescope yields a per-mode rate $r \sim 10^8$\,cps but
$\mathrm{SNR}/\sqrt{T/\mathrm{hr}} \sim 1.5\times10^{-3}$; even
$N \sim 3000$-channel spectrograph multiplexing reaches only ${\sim}7\times10^{-2}$
at $T_{\rm obs} = 10^4$\,hr, far below detection.
\textit{Sirius~B (best-case absorption-line target).} A broad absorption
line can only \emph{reduce} $R$ relative to the unobscured continuum, so
there is no photon-rate advantage from the absorption profile. The deep,
Stark-broadened H$\alpha$ wing of this nearby DA white dwarf
($D = 2.64$\,pc) gives $\mathrm{SNR}/\sqrt{T/\mathrm{hr}} \approx 1.4\times10^{-2}$
at $\alpha = 1/2$ (${\sim}10\times$ below a Mrk~766-class NLS1) and is
entirely negligible at $\alpha = 2/3$. Such targets remain useful for
null-test calibration --- a nearby DA white dwarf has a Lorentzian
profile with negligible predicted foam signal for any plausible $\alpha$.
\textit{The solar Ca~\textsc{ii}\,K line.} Combines proximity with
extended-source dilution ($M \approx 4\times10^{10}$ independent spatial
modes for a 10\,m aperture); reaches
$\mathrm{SNR}/\sqrt{T/\mathrm{hr}} \approx 2\times10^{-6}$, about
$10^4\times$ below the AGN reference.

% ==============================================================
\section{Conclusions}
\label{sec:conclusions}

A single second-order coherence observable of spacetime foam emerges from
this framework: the finite-lag Siegert discrepancy
$\mathcal{S}(\tau) = g^{(2)}(\tau)-1-|g^{(1)}(\tau)|^2$, present when foam
fluctuates on timescales shorter than the optical coherence time
($\tau_{\rm foam}\lesssim\tau_L$).  No foam-induced zero-lag bunching
suppression arises in this model: at $\tau = 0$ the same-direction jitter
increments vanish identically and $g^{(2)}_{\rm obs}(0) = 2$ for any
extended source viewed by a single-mode detector
(Section~\ref{sec:extended}; Appendix~\ref{app:g2ext}).  Its value lies
not in raw sensitivity --- the signal is at ${\sim}10^{-16}$ for the
holographic $\alpha = 2/3$ model, well below current reach --- but in
three structural advantages over existing fringe-visibility and
Strehl-ratio tests:
(i) $\mathcal{S}(\tau)$ is constructed directly from the data, with
$g^{(1)}_{\rm src}(\tau)$ obtained from the source spectrum via
Wiener--Khinchin, so a detection requires no separate source model;
(ii) for an extended source the Siegert discrepancy is attenuated by an
angular-size factor $\Phi_{\rm eff}(r_\perp)\in[0,1]$, providing a
differential probe of the foam transverse coherence length $r_\perp$ at
finite lag --- the point-source value is the maximum-signal upper bound;
(iii) the framework remains sensitive to foam models with long-range
transverse correlations that evade imaging tests
entirely~\cite{LeeZurekChen2024}, and supplies a chromatic discriminant
against plasma scattering through the $\sigma_\phi \propto 1/\lambda$
scaling.

The framework identifies high-priority targets: intrinsically compact
emitters that saturate the upper bound regardless of $r_\perp$
(narrow-line Seyfert~1 cores, lensed-quasar BLRs), and high-flux
transients (Type~IIn supernovae, LBV giant eruptions) whose
photon-rate advantage compensates for the modest cosmological lever arm
at kpc--Mpc distances.  A standing rapid-response HBT plan --- covering
$\eta$\,Car, any galactic core-collapse supernova, and any nearby
($D \lesssim 30$\,Mpc) Type~IIn identified by transient surveys ---
would let the framework exploit such high-yield events without prior
allocation; transient targets with substantial angular extent contribute
to the $r_\perp$ discriminant via $\Phi_{\rm eff}$.  Strongly-lensed
quasars (Section~\ref{sec:candidate_sources}; J1330$-$0905, comparable to
the NLS1 baseline after cosmological correction) round out the persistent
target class.

The achievable sensitivities are derived in Section~\ref{sec:prospects}.
The single finite-lag channel applies to point-like and extended
targets alike, the latter scaled by $\Phi_{\rm eff}$.  What HBT-on-AGN
cannot reach is the holographic prediction $\sigma_\phi \sim 10^{-8}$,
many orders of magnitude below the most aggressive HBT projection and
the natural target of tabletop holographic-noise
interferometers~\cite{VerlindeZurek2021}.  Realising these projections
demands parallel investment in systematic-error control: null-test
calibration sources with negligible predicted foam signal are essential
to verify that instrumental artifacts do not limit sensitivity before
photon statistics do.

This analysis has focused on single-telescope observations, constraining foam models through the temporal statistics
of intensity fluctuations. However, intensity interferometry with spatially separated telescopes on different baselines
offers a complementary and more powerful probe. Cross-correlations between distant detectors encode both the
temporal and spatial structure of spacetime foam, providing sensitivity to foam models whose fluctuations are not
separable into independent temporal and spatial components. Such multi-baseline measurements would break degen-
eracies inherent in single-telescope data and reveal the fundamental spacetime correlations underlying foam-induced
decoherence.

\begin{acknowledgments}
This work was supported in part by the Director, Office of Science, Office of
High Energy Physics of the U.S.\ Department of Energy under Contract No.\
DE-AC02-05CH11231.

The authors used Claude (Anthropic) and Valency Systems to assist with drafting and editing portions of this manuscript along with detailed reference searches. All content was reviewed and verified by the authors, who take full responsibility for the work.

We thank Robin Kaiser for helpful comments.
\end{acknowledgments}

\appendix

\section{Pedestrian derivation of $g^{(2)}_{\rm obs}(\tau)$}
\label{app:g2ext}

This appendix carries out the four-point function calculation underlying
Section~\ref{sec:extended} step by step.  The aim is to expose how the
random variables $(X,Y)$ that enter the connected Wick pairing arise, and
to identify the regime in which the 2D average of Eq.~\eqref{eq:g2ext_correct}
collapses to a 1D convolution (the point-source limit $\Phi=1$) versus
the regime in which the cross-direction carrier-phase factor
$\mathcal{F}_{\rm cross}$ of Eq.~\eqref{eq:Fcross} appears.

\subsection{Setup}

The observed field is Eq.~\eqref{eq:EobsExt},
\begin{equation}
  E_{\rm obs}(t) = \int d\boldsymbol{\theta}\, I^{1/2}(\boldsymbol{\theta})\,
    E_{\boldsymbol{\theta}}\!\big(t - \delta t(t;\boldsymbol{\theta})\big),
  \label{eq:Eobs_app}
\end{equation}
with three statistical inputs:
\begin{enumerate}
  \item Each $E_{\boldsymbol{\theta}}$ is a stationary, zero-mean,
    complex-Gaussian (analytic-signal) process with unit mean intensity
    $\langle |E_{\boldsymbol{\theta}}(s)|^2 \rangle = 1$.
  \item Source spatial incoherence:
    $\langle E^*_{\boldsymbol{\theta}}(s)\,E_{\boldsymbol{\theta}'}(s')\rangle
    = g^{(1)}_{\rm src}(s'-s)\,\delta(\boldsymbol{\theta}-\boldsymbol{\theta}')$,
    with anomalous pairings $\langle EE\rangle$ and $\langle E^*E^*\rangle$
    vanishing.
  \item $\delta t$ is independent of $E$, and its joint distribution at
    distinct $(t,\boldsymbol{\theta})$ is the Gaussian process whose
    covariance is Eq.~\eqref{eq:pathcov} divided by $c^2$.
\end{enumerate}
The brightness distribution is normalized,
$\int I(\boldsymbol{\theta})\,d\boldsymbol{\theta} = 1$, so
$\langle I_{\rm obs}\rangle = 1$ throughout this appendix.

\subsection{Mean intensity}

Plugging Eq.~\eqref{eq:Eobs_app} into $\langle|E_{\rm obs}(t)|^2\rangle$
and using spatial incoherence to collapse the angular integral to the
diagonal,
\begin{align}
  \langle |E_{\rm obs}(t)|^2 \rangle
  &= \int\!\!\int d\boldsymbol{\theta}\,d\boldsymbol{\theta}'\,
    I^{1/2}(\boldsymbol{\theta})I^{1/2}(\boldsymbol{\theta}')\,
    \big\langle E^*_{\boldsymbol{\theta}}(t-\delta t_{1})\,
                E_{\boldsymbol{\theta}'}(t-\delta t_{2})\big\rangle \notag\\
  &= \int d\boldsymbol{\theta}\, I(\boldsymbol{\theta})\,
    \big\langle |E_{\boldsymbol{\theta}}(t-\delta t(t;\boldsymbol{\theta}))|^2 \big\rangle
    = \int d\boldsymbol{\theta}\, I(\boldsymbol{\theta}) = 1,
\end{align}
where $\delta t_1 = \delta t(t;\boldsymbol{\theta})$ and
$\delta t_2 = \delta t(t;\boldsymbol{\theta}')$.  By stationarity, the
random retarded time does not change the marginal second moment of each
$E_{\boldsymbol{\theta}}$, so the foam preserves the mean intensity
exactly: $\langle |E_{\rm obs}|^2 \rangle = \langle |E_{\rm src}|^2\rangle$.

\subsection{Four-point function and Wick contraction}

The intensity autocorrelation is
$\langle I(t)\,I(t+\tau)\rangle = \langle E^*_{\rm obs}(t)\,E_{\rm obs}(t)\,
E^*_{\rm obs}(t+\tau)\,E_{\rm obs}(t+\tau)\rangle$.  Expanding each
$E_{\rm obs}$ as in Eq.~\eqref{eq:Eobs_app} introduces four direction
labels $\boldsymbol{\theta}_1,\ldots,\boldsymbol{\theta}_4$ and four
retarded times
\begin{equation}
\begin{aligned}
  s_1 &= t - \delta t(t;\boldsymbol{\theta}_1), &\quad
  s_2 &= t - \delta t(t;\boldsymbol{\theta}_2), \\
  s_3 &= t+\tau - \delta t(t+\tau;\boldsymbol{\theta}_3), &\quad
  s_4 &= t+\tau - \delta t(t+\tau;\boldsymbol{\theta}_4).
\end{aligned}
\end{equation}
Conditioning on a realization of $\delta t$, the four-point function of
the independent zero-mean complex-Gaussian $E_{\boldsymbol{\theta}}$
obeys Wick's (Isserlis') theorem:
\begin{equation}
  \langle E^*_{\boldsymbol{\theta}_1}(s_1)\,E_{\boldsymbol{\theta}_2}(s_2)\,
    E^*_{\boldsymbol{\theta}_3}(s_3)\,E_{\boldsymbol{\theta}_4}(s_4)\rangle_{E}
  = \langle E^*_1 E_2\rangle\langle E^*_3 E_4\rangle
  + \langle E^*_1 E_4\rangle\langle E^*_3 E_2\rangle.
  \label{eq:Wick_app}
\end{equation}

\paragraph{Disconnected pairing.}
The first term in Eq.~\eqref{eq:Wick_app} forces
$\boldsymbol{\theta}_1=\boldsymbol{\theta}_2$ and
$\boldsymbol{\theta}_3=\boldsymbol{\theta}_4$ via the spatial-incoherence
delta.  With $\boldsymbol{\theta}_1=\boldsymbol{\theta}_2$ the retarded
times coincide ($s_1=s_2$), so the two-point function is
$g^{(1)}_{\rm src}(0)=1$; the same holds for the
$\boldsymbol{\theta}_3,\boldsymbol{\theta}_4$ pair.  The contribution is
$\int I(\boldsymbol{\theta}_1)d\boldsymbol{\theta}_1 \int
I(\boldsymbol{\theta}_3)d\boldsymbol{\theta}_3 = 1 =
\langle I(t)\rangle\langle I(t+\tau)\rangle$, which subtracts cleanly
when forming $g^{(2)}-1$.

\paragraph{Connected pairing.}
The second term forces $\boldsymbol{\theta}_1=\boldsymbol{\theta}_4\equiv
\boldsymbol{\theta}$ and $\boldsymbol{\theta}_3=\boldsymbol{\theta}_2\equiv
\boldsymbol{\theta}'$.  The two surviving two-point functions are
evaluated at retarded-time differences
\begin{equation}
  s_4 - s_1 = \tau - X,\qquad
  s_2 - s_3 = -(\tau - Y),
\end{equation}
with the two \emph{same-direction temporal increments}
\begin{equation}
  X \equiv \delta t(t+\tau;\boldsymbol{\theta}) - \delta t(t;\boldsymbol{\theta}),
  \qquad
  Y \equiv \delta t(t+\tau;\boldsymbol{\theta}') - \delta t(t;\boldsymbol{\theta}')
\end{equation}
[Eq.~\eqref{eq:XYdef} of the body].
Using $g^{(1)}_{\rm src}(-u) = g^{(1)*}_{\rm src}(u)$, the connected
contribution to $g^{(2)}_{\rm obs}-1$ is
\begin{equation*}
  \boxed{\;
  g^{(2)}_{\rm obs}(\tau) - 1
    = \int\!\!\int d\boldsymbol{\theta}\,d\boldsymbol{\theta}'\;
      I(\boldsymbol{\theta})I(\boldsymbol{\theta}')\;
      \big\langle g^{(1)}_{\rm src}(\tau - X)\,
                 g^{(1)*}_{\rm src}(\tau - Y)\big\rangle_{X,Y}.
  \;}
\end{equation*}
This is Eq.~\eqref{eq:g2ext_correct} of the body.

\subsection{Joint distribution of $(X,Y)$}

$X$ and $Y$ are linear functionals of the Gaussian process $\delta t$,
hence jointly Gaussian and zero-mean.  Using
Eq.~\eqref{eq:pathcov} and $C_\alpha$ even in $\tau$,
\begin{align}
  \operatorname{Var}(X) = \operatorname{Var}(Y)
    &= \frac{2\sigma_\ell^2(D)}{c^2}\big[1 - C_\alpha(\tau;D)\big],
    \label{eq:varX}\\
  \operatorname{Cov}(X,Y)
    &= \frac{2\sigma_\ell^2(D)}{c^2}\big[1 - C_\alpha(\tau;D)\big]\,
       \Phi\!\left(b/r_\perp\right),
    \label{eq:covXY}
\end{align}
so the correlation coefficient is $\rho = \Phi(b/r_\perp)$.

The structural point: $X$ is the temporal jitter increment
\emph{along direction $\boldsymbol{\theta}$}, and $Y$ is the increment
\emph{along direction $\boldsymbol{\theta}'$}.  Neither involves a
cross-direction sample of $\delta t$; the transverse coherence factor
$\Phi$ enters only through their correlation.  Equation~\eqref{eq:g2ext_correct}
is therefore a 2D average over the joint distribution of $(X,Y)$, not a
1D convolution with a single variable $\Delta$.

It collapses to a 1D convolution only when $\rho = 1$, i.e.\ when
$\Phi = 1$ (point source, or the transversely-coherent limit
$b\ll r_\perp$).  In that case $X = Y$ almost surely,
$g^{(1)}_{\rm src}(\tau-X)g^{(1)*}_{\rm src}(\tau-X) =
|g^{(1)}_{\rm src}(\tau-X)|^2 = g^{(2)}_{\rm src}(\tau-X)-1$, and the
point-source formula Eq.~\eqref{eq:g2obs} is recovered with
$\sigma_\Delta^2 = \operatorname{Var}(X) =
(2\sigma_\ell^2/c^2)[1-C_\alpha(\tau)]$.

\subsection{Zero lag: $g^{(2)}_{\rm obs}(0) = 2$ exactly}

At $\tau=0$ each of the increments in Eq.~\eqref{eq:XYdef} is the
difference of a random variable with itself:
\begin{equation}
  X\big|_{\tau=0} = 0,\qquad Y\big|_{\tau=0} = 0,
\end{equation}
for every realization of $\delta t$ and every pair
$(\boldsymbol{\theta},\boldsymbol{\theta}')$.  This is independent of
$\Phi$, $C_\alpha$, $\sigma_\ell$, and $D$.  Equation~\eqref{eq:g2ext_correct} reduces to
\begin{equation}
  g^{(2)}_{\rm obs}(0) - 1
  = \int\!\!\int d\boldsymbol{\theta}\,d\boldsymbol{\theta}'\;
    I(\boldsymbol{\theta})I(\boldsymbol{\theta}')\,
    |g^{(1)}_{\rm src}(0)|^2 = 1,
\end{equation}
so $g^{(2)}_{\rm obs}(0) = 2$ exactly.

The same conclusion follows from a direct argument that bypasses the
Wick contraction.  Conditioning on $\delta t$,
$E_{\rm obs}(t) = \int d\boldsymbol{\theta}\,I^{1/2}(\boldsymbol{\theta})\,
\widetilde E_{\boldsymbol{\theta}}$ with
$\widetilde E_{\boldsymbol{\theta}} \equiv
E_{\boldsymbol{\theta}}(t-\delta t(t;\boldsymbol{\theta}))$.  By
stationarity each $\widetilde E_{\boldsymbol{\theta}}$ has the same
distribution as $E_{\boldsymbol{\theta}}(t)$; since the $E_{\boldsymbol{\theta}}$
are independent across $\boldsymbol{\theta}$, so are the
$\widetilde E_{\boldsymbol{\theta}}$.  Hence $E_{\rm obs}(t)$ is a linear
combination of independent zero-mean complex Gaussians with fixed
weights $I^{1/2}(\boldsymbol{\theta})$, and is itself a zero-mean complex
Gaussian.  The Gaussian-moment identity gives
\begin{equation}
  \big\langle |E_{\rm obs}(t)|^4 \,\big|\, \delta t \big\rangle
  = 2\,\big\langle |E_{\rm obs}(t)|^2 \,\big|\, \delta t \big\rangle^2 = 2
\end{equation}
for every $\delta t$, so $g^{(2)}_{\rm obs}(0) = \langle I^2\rangle/\langle
I\rangle^2 = 2$ unconditionally.  Foam imprints itself on multi-time
correlations of $E_{\rm obs}$ but not on its single-time amplitude
statistics; the standard Gaussian identity $g^{(2)}(0)=2$ for any
zero-mean complex Gaussian field is preserved.

\subsection{Finite lag: narrow-band reduction}

For $\tau\neq 0$, split $g^{(1)}_{\rm src}$ into carrier and envelope
(Section~\ref{sec:narrowband}),
$g^{(1)}_{\rm src}(\tau) = e^{i\omega_0\tau}\,V_{\rm src}(\tau)$ with
$V_{\rm src}$ real and slowly varying on the carrier scale, so
\begin{equation}
  g^{(1)}_{\rm src}(\tau-X)\,g^{(1)*}_{\rm src}(\tau-Y)
  = e^{-i\omega_0(X-Y)}\,V_{\rm src}(\tau-X)\,V_{\rm src}(\tau-Y).
\end{equation}
Change variables in the joint Gaussian average to the sum and difference
\begin{equation}
  U \equiv \tfrac{1}{2}(X+Y),\qquad V \equiv X - Y,
\end{equation}
which (because $X,Y$ have the same variance) are independent zero-mean
Gaussians with
\begin{equation}
  \sigma_U^2 = \frac{\sigma_\ell^2(D)}{c^2}\big[1-C_\alpha(\tau;D)\big](1+\Phi),
  \qquad
  \sigma_V^2 = \frac{4\sigma_\ell^2(D)}{c^2}\big[1-C_\alpha(\tau;D)\big](1-\Phi).
  \label{eq:sigmaUV}
\end{equation}
In the narrow-band regime $\sigma_V\ll\tau_L$, so
$V_{\rm src}(\tau-U-V/2)\,V_{\rm src}(\tau-U+V/2) \approx
V_{\rm src}^2(\tau-U)$, and the $V$ integral is the characteristic
function of a Gaussian:
\begin{equation}
  \int dV\,\mathcal{N}(V;0,\sigma_V^2)\,e^{-i\omega_0 V}
  = e^{-\omega_0^2\sigma_V^2/2}
  = \exp\!\Big(-2\sigma_\phi^2(D)\,[1-C_\alpha(\tau;D)]\,[1-\Phi]\Big),
\end{equation}
with $\sigma_\phi = \omega_0\sigma_\ell/c$ the rms accumulated carrier
phase along a single ray.  The result reproduces Eq.~\eqref{eq:g2ext_corrected}
of the body, with the cross-direction carrier-phase factor
$\mathcal{F}_{\rm cross}(\tau;b)$ of Eq.~\eqref{eq:Fcross}.

\subsection{Summary}

Equations~\eqref{eq:g2ext_correct} and \eqref{eq:g2ext_corrected} of the
body are the key results of this pedestrian derivation.  Two structural
features deserve emphasis: a cross-direction carrier-phase factor
$\mathcal{F}_{\rm cross}$ controlled by the imaging-channel scale
$\sigma_\phi$ appears (so the $\sigma_\phi$/$\sigma_\Delta/\tau_c$
decoupling of Section~\ref{sec:scales} is exact only for a point source);
and the envelope-convolution variance is $\sigma_U^2 =
(\sigma_\ell^2/c^2)[1-C_\alpha](1+\Phi)$ rather than the point-source
$\sigma_\Delta^2 = (2\sigma_\ell^2/c^2)[1-C_\alpha]$, the two coinciding
only at $\Phi=1$.  At $\tau=0$ both $\mathcal{F}_{\rm cross}$ and
$\sigma_U^2$ trivialize ($C_\alpha(0)=1$), reproducing the rigorous
single-time identity $g^{(2)}_{\rm obs}(0) = 2$ for any
$\Phi_{\rm eff}\in[0,1]$.

\bibliographystyle{apsrev4-2}
\bibliography{references}

\end{document}
